% Mathematics appendix, Astra, 2026-09-10.
% Integration: input AFTER \appendix. Requires the macros/environments and
% amsmath,amssymb,amsthm,mathtools,natbib of paper/main.tex.
% All local labels use astra_ to avoid collisions. Main-text cor:attribution
% should state the sharp result proved in Corollary cor:astra_attribution.
% Do not input the old and new proofs as duplicate theorem statements.
% BibTeX records for every citation used here are in the final comment block.
% Review-note numbers (1.1, 3.2, N, etc.) are mnemonic names, not hard-coded
% manuscript counters. The existing main-text labels remain unchanged.

\section{Mathematical scope and normalization}
\label{sec:astra_math}
We use the parameter blocks $\thetap,\phip$ of the main text and
$z=N^{-1/2}\operatorname{stack}(\hat y_n)$, $r=\nabla_z\ell(z)$,
$J_b=\partial z/\partial b$, $\loss=\ell\circ z$. Write
\[
 K_b=J_bJ_b^\top,\qquad G_{bb}=J_b^\top H_\tau J_b,
 \qquad \nabla_b\loss=J_b^\top r.
\]
Thus $K_b$ is an output-space logit kernel; it is not the parameter-space
GGN block $G_{bb}$. For averaged softmax CE,
$r=N^{-1/2}\operatorname{stack}(p_n-y_n)$ and $H_\tau$ is the block-diagonal
softmax Hessian. For normalized squared loss, $r=z-y$ and $H_\tau=I$.
All parameter norms are Euclidean (Frobenius for matrix parameters);
matrix-map bounds are operator norms unless explicitly marked otherwise.
Here $C=N_{\mathrm{cls}}$ denotes output dimension; envelope amplitudes
are written $C_r$ to avoid a collision. A normalized convention removes
extraneous powers of $N$, but does not make data-dependent constants
such as kernel minimum eigenvalues independent of $N$.

Theorem~\ref{thm:hessian} and Lemma~\ref{lem:phi_scaling} are geometric
statements. Theorem~\ref{thm:twoscale} additionally requires a trajectory
Jacobian cap and a residual envelope. Neither a large
$\lambda_{\max}(G_{\phip\phip})$ nor parameter count alone supplies
residual-subspace coercivity or a comparison with a separately trained
frozen model. For unit gradient flow, $\dot z=-(K_\thetap+K_\phip)r$;
under CE, $\dot r=-H_\tau(K_\thetap+K_\phip)r$, not the squared-loss
constant-kernel residual equation.

\section{Discrete displacement, clipping, and optimizer conventions}
\label{sec:astra_discrete}
The following seven statements collect the bounds numbered (4.1)--(4.7)
in the review notes and distinguish the two clipping algorithms. They are
pathwise: no convergence, stochastic independence, or smoothness assumption
is needed beyond the stated gradient bound at every iterate. There is no
encoder weight decay unless explicitly included. A uniform bound
$\|g_k\|\le B_kR_k\le BR_k$ derived from the main theorem requires its
containment/Jacobian hypothesis at \emph{each discrete iterate}. A full-data
or smoothed residual envelope does not establish a sampled-gradient envelope.
An observed residual floor limits a fitted envelope; a positive floor is
not a universal property of mini-batch training.

\begin{lemma}[Pathwise GD and clipping before a GD update]
\label{lem:astra_gd}
Let $K\ge0$, $\eta_k\ge0$, $a_k\in[0,1]$, and
$\|g_k\|\le B_kR_k$, where $B_k,R_k\ge0$. If
$\thetap_{k+1}=\thetap_k-\eta_ka_kg_k$, then
\begin{equation}\label{eq:astra_gd}
 \|\thetap_K-\thetap_0\|
 \le\sum_{k=0}^{K-1}\eta_ka_kB_kR_k
 \le\sum_{k=0}^{K-1}\eta_kB_kR_k.
\end{equation}
The clipping scalar may depend on all parameter groups and the history.
\end{lemma}
\begin{proof}
Telescope the updates and apply the triangle inequality.
\end{proof}

\begin{corollary}[Polynomial envelope in step units]
\label{cor:astra_steps}
In Lemma~\ref{lem:astra_gd}, suppose $\eta_k=\eta\ge0$, $B_k\le B$,
and $R_k\le C_r/(1+\mu_s k)$, with $C_r,\mu_s>0$.
For $K\ge1$,
\begin{equation}\label{eq:astra_steps}
 \|\thetap_K-\thetap_0\|
 \le\eta BC_r\left[1+\frac{\log(1+\mu_s(K-1))}{\mu_s}\right].
\end{equation}
For $0<\delta<C_r$, suppose the envelope holds through
$K_* =\lceil(C_r/\delta-1)/\mu_s\rceil$. Then the first index
$K_\delta=\min\{k:R_k\le\delta\}$ satisfies $K_\delta\le K_*$ and
\begin{equation}\label{eq:astra_hit}
 \|\thetap_{K_\delta}-\thetap_0\|
 \le\eta BC_r[1+\mu_s^{-1}\log(C_r/\delta)].
\end{equation}
The displacement is zero if $K_\delta=0$.
\end{corollary}
\begin{proof}
For decreasing $f(x)=(1+\mu_sx)^{-1}$,
$\sum_{k=1}^{K-1}f(k)\le\int_0^{K-1}f(x)\,dx$.
The envelope gives $R_{K_*}\le\delta$. If $K_\delta\ge1$,
$K_\delta-1\le K_*-1<(C_r/\delta-1)/\mu_s$; substitute in
\eqref{eq:astra_steps}.
\end{proof}
\begin{remark}[Flow-to-step dictionary]
For constant positive $\eta$, $t_k=\eta k$ and $\mu_s=\eta\mu_f$.
Consequently $\eta BC_r/\mu_s=BC_r/\mu_f$. Omitting the leading
$\eta$ in \eqref{eq:astra_steps} is a units error. The extra
$\eta BC_r$ is a left-sum allowance, not an additional mechanism.
\end{remark}

\begin{lemma}[Variable-step envelope]
\label{lem:astra_variable}
Under Lemma~\ref{lem:astra_gd}, let $t_k=\sum_{i<k}\eta_i$,
$B_k\le B$, $R_k\le C_r/(1+\mu_ft_k)$ with $C_r,\mu_f>0$.
For $K\ge1$ and $\eta_{\max}=\max_{k<K}\eta_k$,
\begin{equation}\label{eq:astra_variable}
 \|\thetap_K-\thetap_0\|\le BC_r\left[
 \frac{\log(1+\mu_ft_K)}{\mu_f}
 +\eta_{\max}\left(1-\frac1{1+\mu_ft_K}\right)\right].
\end{equation}
\end{lemma}
\begin{proof}
The excess of the left sum for $f(t)=(1+\mu_ft)^{-1}$ over its
integral on $[t_k,t_{k+1}]$ is at most
$\eta_k(f(t_k)-f(t_{k+1}))$. Bound $\eta_k$ by $\eta_{\max}$ and
telescope. Zero step sizes cause no difficulty.
\end{proof}

\begin{lemma}[Momentum, including gradient clipping before accumulation]
\label{lem:astra_momentum}
Let $0\le\beta<1$, $\eta_k\ge0$, $a_k\in[0,1]$, and
$\|g_k\|\le B_kR_k$. For
$v_{k+1}=\beta v_k+a_kg_k$ and
$\thetap_{k+1}=\thetap_k-\eta_kv_{k+1}$,
\begin{align}
 \thetap_K-\thetap_0
 &=-v_0\sum_{k<K}\eta_k\beta^{k+1}
   -\sum_{i<K}a_ig_i\sum_{k=i}^{K-1}\eta_k\beta^{k-i}.
 \label{eq:astra_momentum_exact}
\end{align}
For constant $\eta$ this implies
\begin{equation}\label{eq:astra_momentum}
 \|\thetap_K-\thetap_0\|\le
 \frac{\eta\beta(1-\beta^K)}{1-\beta}\|v_0\|
 +\eta\sum_{i<K}\frac{1-\beta^{K-i}}{1-\beta}a_iB_iR_i.
\end{equation}
For $v_0=0$, \eqref{eq:astra_steps}--\eqref{eq:astra_hit} therefore
remain valid with an additional factor $1/(1-\beta)$.
\end{lemma}
\begin{proof}
Induction gives $v_{k+1}=\beta^{k+1}v_0+
\sum_{i=0}^k\beta^{k-i}a_ig_i$. Substitute into the parameter update,
interchange the finite sums, and sum the geometric series. For variable
steps the inner coefficient is at most $\eta_{\max}/(1-\beta)$.
\end{proof}
\begin{remark}
This is zero-dampening, non-Nesterov momentum. It includes global
\emph{gradient} clipping through its actual $a_k$; the gradients and
later iterates belong to the clipped run itself. It does not compare
a clipped run's displacement with a separately evolved unclipped run.
The factor $1/(1-\beta)$ is approached by long sequences of aligned gradients.
\end{remark}

\begin{lemma}[Velocity clipping: sharper norm induction]
\label{lem:astra_velocity}
Let $c_k\ge0$, $0\le\beta<1$, and let
$\operatorname{clip}_{c}(u)=u\min\{1,c/\|u\|\}$, with
$\operatorname{clip}_{c}(0)=0$. Suppose
\[
 v_{k+1}=\operatorname{clip}_{c_k}(\beta v_k+g_k),\qquad
 \thetap_{k+1}=\thetap_k-\eta_kv_{k+1},\quad \eta_k\ge0.
\]
Define the \emph{comparison} coefficient
$a_k=\min\{1,c_k/\|g_k\|\}$ for $g_k\ne0$ and $a_k=1$ otherwise.
Then
\[
 \|v_{k+1}\|\le\beta^{k+1}\|v_0\|
       +\sum_{i=0}^k\beta^{k-i}a_i\|g_i\|.
\]
Consequently \eqref{eq:astra_momentum} holds when
$\|g_i\|\le B_iR_i$. For variable steps the corresponding norm bound
uses the same finite weights as \eqref{eq:astra_momentum_exact}.
The vector identity \eqref{eq:astra_momentum_exact} is not asserted.
\end{lemma}
\begin{proof}
Write $x=\beta\|v_k\|$, $y=\|g_k\|$. Then
\[
 \|v_{k+1}\|\le\min(c_k,x+y)\le x+\min(c_k,y)
 =\beta\|v_k\|+a_k\|g_k\|.
\]
If $y\le c_k$ use the second argument of the minimum; if $y>c_k$ use
its first. Induct and sum the parameter-step norms.
\end{proof}
\begin{remark}[Two valid coefficient conventions]
The comparison $a_k$ above differs from the actual multiplier
$t_k=\min\{1,c_k/\|\beta v_k+g_k\|\}$, defined as one at a zero
argument. The same norm bound is also valid with $t_i$ in place of
$a_i$, by the separate induction
$\|v_{k+1}\|\le\beta\|v_k\|+t_k\|g_k\|$.
In fact the exact unrolling is
\[
 v_{k+1}=\beta^{k+1}\left(\prod_{j=0}^kt_j\right)v_0
 +\sum_{i=0}^k\beta^{k-i}\left(\prod_{j=i}^kt_j\right)g_i.
\]
Dropping subsequent multipliers gives the claimed $t_i$ bound.
The two coefficient conventions need not be ordered, owing to gradient
cancellation. State which one is used. For radial implementations with
$c/(\|u\|+\varepsilon_{\rm clip})$ and clipping multiplier in $[0,1]$,
this actual-multiplier argument still applies; the ideal-clipping
comparison $a_i$ is not obtained by silently equating the two formulas.
\end{remark}

\begin{lemma}[Positive-stabilizer Adam]
\label{lem:astra_adam}
Suppose the first moment is initialized at zero, $0\le\beta_1<1$,
$0\le\eta_k\le\eta_{\max}$, and
\[
 \widehat m_k=\sum_{i=0}^k
 \frac{(1-\beta_1)\beta_1^{k-i}}{1-\beta_1^{k+1}}g_i,
 \qquad \thetap_{k+1}-\thetap_k=-\eta_kD_k\widehat m_k,
 \qquad \|D_k\|\le\zeta^{-1},\quad\zeta>0.
\]
Then
\begin{equation}\label{eq:astra_adam}
 \|\thetap_K-\thetap_0\|
 \le\frac{\eta_{\max}}{\zeta(1-\beta_1)}\sum_{i<K}\|g_i\|.
\end{equation}
Standard Adam denominators $\sqrt{\widehat v_k}+\zeta$ satisfy the
preconditioner bound when $\widehat v_k\ge0$ coordinatewise. If
clipping occurs before moment accumulation, $g_i$ here is the gradient
\emph{after} that clipping.
\end{lemma}
\begin{proof}
Bound the update by $\eta_{\max}\zeta^{-1}\|\widehat m_k\|$ and
interchange sums. For each $i$, the sum of its nonnegative weights is
at most $\sum_{k\ge i}\beta_1^{k-i}\le(1-\beta_1)^{-1}$, since
$(1-\beta_1)/(1-\beta_1^{k+1})\le1$.
\end{proof}
\begin{remark}[Numerically weak optimizer constant]
The factor $1/\zeta$ makes this certificate numerically vacuous at
standard small stabilizers unless a much smaller preconditioner bound
is independently certified. At $\zeta=0$, the first scalar Adam step
has magnitude $\eta$ for every nonzero $g_0$, so no constant uniform
in residual scale bounds it by a multiple of $\eta\|g_0\|$.
\end{remark}

\begin{lemma}[Decoupled encoder decay must be included]
\label{lem:astra_decay}
If an update is $\thetap_{k+1}-\thetap_k=-u_k-\eta_k\lambda_k\thetap_k$
with $\eta_k,\lambda_k\ge0$, then
$\|\thetap_K-\thetap_0\|\le\sum_{k<K}\|u_k\|+
\sum_{k<K}\eta_k\lambda_k\|\thetap_k\|$.
Even with $u_k=0$, constant decay $\lambda$ gives
$\thetap_K=\prod_{k<K}(1-\eta_k\lambda)\thetap_0$.
\end{lemma}
\begin{proof}
Telescope for the bound and iterate the scalar multiplication for the
identity. A nonzero decay displacement can occur at identically zero
residual, so a residual-only bound cannot omit this term.
\end{proof}

\section{Joint and frozen trajectory comparisons}
\label{sec:astra_bridges}
\begin{theorem}[Monotone frozen flow; review Theorem 1.1]
\label{thm:astra_monotone}
On $[0,T]$ let the joint and frozen paths be absolutely continuous,
with common spatial initialization and
$\dot\phip_J=-\nabla_\phip\loss(\thetap_J,\phip_J)$,
$\dot\phip_F=-\nabla_\phip\loss(\thetap_0,\phip_F)$ a.e.
Let $d\ge0$ be locally integrable with
$\|\thetap_J(t)-\thetap_0\|\le d(t)$.
Assume the differentiable frozen objective has an $\alpha$-monotone
gradient on a region containing both spatial paths, $\alpha\ge0$:
\[
 \langle v-w,\nabla_\phip\loss(\thetap_0,v)
 -\nabla_\phip\loss(\thetap_0,w)\rangle\ge\alpha\|v-w\|^2.
\]
Assume finite nonnegative constants $K_{\theta\phi},L_\theta,L_\phi$
satisfy, on the parameter pairs used below,
\begin{align*}
 \|\nabla_\phip\loss(\thetap_J,\phip_J)
       -\nabla_\phip\loss(\thetap_0,\phip_J)\|&\le K_{\theta\phi}d(t),\\
 \|z(\thetap_J,\phip_J)-z(\thetap_0,\phip_J)\|&\le L_\theta d(t),\\
 \|z(\thetap_0,\phip_J)-z(\thetap_0,\phip_F)\|&\le L_\phi\|\phip_J-\phip_F\|.
\end{align*}
Then
\begin{align}
 \|\phip_J(t)-\phip_F(t)\|&\le K_{\theta\phi}\int_0^t e^{-\alpha(t-s)}d(s)\,ds,
 \label{eq:astra_monotone}\\
 \|z_J(t)-z_F(t)\|&\le L_\theta d(t)+
 L_\phi K_{\theta\phi}\int_0^t e^{-\alpha(t-s)}d(s)\,ds.
\end{align}
If $d$ is nondecreasing, the last bound is
$[L_\theta+L_\phi K_{\theta\phi}\psi_\alpha(t)]d(t)$, where
$\psi_\alpha(t)=(1-e^{-\alpha t})/\alpha$ for $\alpha>0$ and
$\psi_0(t)=t$.
\end{theorem}
\begin{proof}
For $\Delta=\phip_J-\phip_F$, split the derivative into a frozen-gradient
difference and an encoder-induced forcing. Then
$\frac12(d/dt)\|\Delta\|^2\le-\alpha\|\Delta\|^2+
K_{\theta\phi}d\|\Delta\|$. The scalar norm comparison (or the
regularization $\sqrt{\|\Delta\|^2+\epsilon^2}$ followed by
$\epsilon\downarrow0$) proves \eqref{eq:astra_monotone}. Add and subtract
$z(\thetap_0,\phip_J)$ for the output bound; monotonicity of $d$ gives
the final expression.
\end{proof}
\begin{remark}
Under Theorem~\ref{thm:twoscale}, one may use
$d(t)=BC_r\log(1+\mu t)/\mu$ on its valid horizon. If the first fitting
time lies in that horizon, replace this logarithm by $\log(C_r/\delta)$
only up to that fitting time. The constants $L_\theta,L_\phi,K_{\theta\phi}$
and $\alpha$ are additional path-dependent hypotheses; all must be
controlled across widths before inferring a width rate. A PL inequality
alone supplies none of this monotonicity.
\end{remark}

\begin{proposition}[Counterexample 1.2: PL does not imply contraction]
\label{prop:astra_pl_counter}
For $a,b>0$, $f(x,y)=\frac12(x+ay^2)^2$ satisfies
$\frac12\|\nabla f\|^2\ge f-\inf f$ globally, but nearby gradient-flow
trajectories can expand by an arbitrarily large factor at time one.
\end{proposition}
\begin{proof}
$\inf f=0$ and $\|\nabla f\|^2=(x+ay^2)^2(1+4a^2y^2)\ge2f$.
Along $(x(t),y(t))=(-be^{-t},0)$, transverse linearized variation satisfies
$\dot\xi=2abe^{-t}\xi$, hence
$\xi(t)=\xi(0)\exp(2ab(1-e^{-t}))$. Vary $a$ with the PL constant fixed.
This uses the PL convention of \citet{astra_karimi2016}.
\end{proof}

\begin{theorem}[Affine squared-loss bridge; review Theorem 1.3]
\label{thm:astra_affine}
Let $z=z_0+A(\thetap-\thetap_0)+B(\phip-\phip_0)$ and
$\loss=\frac12\|z-y\|^2$. Train the joint model and the model with
$\thetap=\thetap_0$ by unit Euclidean gradient flow from the same
initialization. Let $K_\thetap=AA^\top$, $K_\phip=BB^\top$.
Suppose a common invariant linear subspace $\mathcal V$ contains
$e_0=z_0-y$ and, on $\mathcal V$,
$K_\phip\succeq\kappa I$, $0\preceq K_\thetap\preceq aI$ with
$\kappa>0$, $a\ge0$. For $t\ge0$,
\begin{align}
 \|z_J(t)-z_F(t)\|&\le at e^{-\kappa t}\|e_0\|,
 &\sup_{t\ge0}\|z_J(t)-z_F(t)\|&\le\frac{a\|e_0\|}{e\kappa},
 \label{eq:astra_affine_gap}\\
 \|\thetap_J(t)-\thetap_0\|&\le
 \frac{\sqrt a}{\kappa}(1-e^{-\kappa t})\|e_0\|,
 &\int_0^\infty\|\nabla_\thetap\loss\|^2dt&\le
 \frac{a\|e_0\|^2}{2\kappa}.\label{eq:astra_affine_disp}
\end{align}
At every finite horizon with positive loss decrease the encoder share
is at most $a/(a+\kappa)$, by Corollary~\ref{cor:astra_attribution}.
\end{theorem}
\begin{proof}
On $\mathcal V$, both semigroups have norm at most $e^{-\kappa t}$.
Their difference obeys the noncommutative Duhamel identity
\[
 e_J(t)-e_F(t)=-\int_0^t e^{-K_\phip(t-s)}K_\thetap
                e^{-(K_\phip+K_\thetap)s}e_0\,ds.
\]
This gives $at e^{-\kappa t}\|e_0\|$, maximized at $t=1/\kappa$.
Since $\|A^\top v\|\le\sqrt a\|v\|$ for $v\in\mathcal V$,
integrate $\dot\thetap=-A^\top e_J$ and its squared norm to get
\eqref{eq:astra_affine_disp}. Both residuals converge to zero.
\end{proof}

\begin{corollary}[Sharp finite-horizon attribution; main Corollary~\ref{cor:attribution}]
\label{cor:astra_attribution}
Explicitly: $\loss=\ell\circ z$ is differentiable, the parameter path
is absolutely continuous, the residual is $r=\nabla_z\ell$ with
$K_\thetap=J_\thetap J_\thetap^\top$,
$K_\phip=J_\phip J_\phip^\top$, both blocks follow unit Euclidean
gradient flow, and the chain rule holds a.e. Suppose for constants
$a\ge0$, $\kappa>0$ that the Rayleigh inequalities
\[
 r^\top K_\thetap r\le a\|r\|^2,\qquad
 r^\top K_\phip r\ge\kappa\|r\|^2
\]
hold at almost every time along the realized trajectory on $[0,T]$,
and $\loss(0)-\loss(T)>0$. Then
\begin{equation}\label{eq:astra_attribution}
 \frac{\int_0^T\|\nabla_\thetap\loss\|^2dt}{\loss(0)-\loss(T)}
 \le\frac{a}{a+\kappa}.
\end{equation}
No common invariant subspace is required: that hypothesis of
Theorem~\ref{thm:astra_affine} is a device for guaranteeing the
Rayleigh bounds along the path, and here they are assumed directly.
\end{corollary}
\begin{proof}
Put $q_b=r^\top K_b r=\|\nabla_b\loss\|^2$.
Pointwise $\kappa q_\thetap\le a q_\phip$. Integrate and use
$\loss(0)-\loss(T)=\int_0^T(q_\thetap+q_\phip)dt$.
Scalar kernels $K_\thetap=a$, $K_\phip=\kappa$ attain equality.
\end{proof}
\begin{remark}
The same result holds for the nonsmooth flow convention defined below
whenever its a.e. energy identity and Rayleigh inequalities hold.
The unit rate is not cosmetic: with per-block rates $\eta_\theta\ne
\eta_\phi$, the numerator $\int\|\nabla_\thetap\loss\|^2dt$ is no
longer the encoder share. For positive constant rates, the actual share
uses $\eta_\theta$ in the numerator and is bounded by
$\eta_\theta a/(\eta_\theta a+\eta_\phi\kappa)$.
\end{remark}

\begin{proposition}[Counterexample 1.4: top eigenvalues are insufficient]
\label{prop:astra_top_counter}
Let $K_\phip=\operatorname{diag}(M,0)$,
$K_\thetap=\operatorname{diag}(0,1)$, $e_0=(0,1)^\top$, $M>0$.
In the affine squared-loss model, the joint/frozen logit gap is
$1-e^{-t}$ although $\lambda_{\max}(K_\phip)=M$.
\end{proposition}
\begin{proof}
$e_F(t)=(0,1)^\top$ and $e_J(t)=(0,e^{-t})^\top$. Each kernel is
realized by its PSD square root as an affine Jacobian. The fast spatial
direction contains none of the residual.
\end{proof}

\begin{proposition}[Robustness of the affine comparison]
\label{prop:astra_robust}
Under the kernel/subspace assumptions of Theorem~\ref{thm:astra_affine},
suppose instead the residuals solve
$\dot e_J=-(K_\phip+K_\thetap)e_J+\xi_J$ and
$\dot e_F=-K_\phip e_F+\xi_F$, with common $e_0$ and locally
integrable $\xi_J,\xi_F\in\mathcal V$. Then add
\[
 \int_0^t e^{-\kappa(t-s)}(\|\xi_J(s)\|+\|\xi_F(s)\|)\,ds
\]
to the first bound in \eqref{eq:astra_affine_gap}. If
$\|\xi_J\|\le\nu_J$ and $\|\xi_F\|\le\nu_F$, the additional
uniform term is at most $(\nu_J+\nu_F)/\kappa$.
\end{proposition}
\begin{proof}
Compare each perturbed solution with its own constant-kernel unperturbed
solution by variation of constants, then use the triangle inequality.
\end{proof}

\begin{proposition}[Logit margins imply a disagreement bound]
\label{prop:astra_margins}
If normalized logits satisfy $\|z_J-z_F\|\le h$ and at most a fraction
$b_\gamma$ of the N samples have frozen top-two unnormalized logit
margin at most $\gamma>0$, the disagreement fraction is at most
$b_\gamma+2h^2/\gamma^2$ (and at most one).
\end{proposition}
\begin{proof}
On a remaining disagreement, for the frozen winner $c_*$ and the joint
winner $c$, the logit perturbation $d_n$ satisfies
$d_{n,c}-d_{n,c_*}>\gamma$. Thus
$\|d_n\|^2\ge(d_{n,c}-d_{n,c_*})^2/2>\gamma^2/2$.
Average and apply Markov. Macro-F1 needs further class-specific
confusion and denominator control.
\end{proof}
\subsection{A finite-data initialization certificate}
\begin{corollary}[Biased-ReLU initialization; review Corollary 1.5]
\label{cor:astra_init}
Fix $N<\infty$, $x_n\in\R^S$, and a deterministic encoder matrix
$\Theta_0$ for which $u_n=\Theta_0x_n\in\R^K$ are pairwise distinct
and $\|u_n\|\le R$. Let
\[
 h_i(u)=\operatorname{ReLU}(w_i^\top u+b_i),\qquad \hat y=Vh+b_2,
\]
with independent $w_i\sim\mathcal N(0,s_w^2I_K/K)$,
$b_i\sim\mathcal N(0,s_b^2)$, $V_{ci}\sim\mathcal N(0,s_w^2/M)$,
$b_{2,c}\sim\mathcal N(0,s_b^2)$, and $s_w,s_b>0$, $1\le c\le C$.
Let $A,B$ be the normalized encoder/spatial logit Jacobians at this
initialization, including all hidden and readout parameters in the
spatial block. Train the resulting \emph{exactly affine} model with
squared loss as in Theorem~\ref{thm:astra_affine}.
Define
\[
 \psi=(h_1(u_n))_{n=1}^N,\quad K_*=N^{-1}\E\psi\psi^\top,
 \quad\kappa_*=\lambda_{\min}(K_*),\quad
 s_{\max}^2=s_w^2R^2/K+s_b^2.
\]
Then $\kappa_*>0$. For $0<\rho<1$, put
\begin{align*}
 M_0&=\left\lceil\frac{18s_{\max}^4}{\rho\kappa_*^2}\right\rceil,
 & A_\rho&=\frac{3C_G}{\rho},
 & R_\rho&=\sqrt{3C_r/\rho},\\
 C_G&=Cs_w^4\operatorname{tr}(\Sigma_X),
 & C_r&=2C(s_w^2s_{\max}^2/2+s_b^2)+2\|y\|^2,
 &\Sigma_X&=N^{-1}\sum_nx_nx_n^\top.
\end{align*}
For each $M\ge M_0$, with probability at least $1-\rho$,
\[
 K_\phip\succeq\kappa_*M I/2,\qquad \|K_\thetap\|\le A_\rho,
 \qquad\|e_0\|\le R_\rho.
\]
Consequently
\[
 \sup_t\|z_J-z_F\|\le\frac{2A_\rho R_\rho}{e\kappa_*M},
 \qquad\sup_t\|\thetap_J-\thetap_0\|
 \le\frac{2\sqrt{A_\rho}R_\rho}{\kappa_*M}.
\]
The displayed sufficient threshold necessarily satisfies
$M_0\ge72N^2/\rho$; it is an existence certificate, not a practical
onset width or a lower bound on widths at which coercivity actually holds.
\end{corollary}
\begin{proof}
If $c^\top K_*c=0$, the continuous function
$\sum_nc_n\operatorname{ReLU}(w^\top u_n+b)$ vanishes almost surely
under a Gaussian with full support, hence everywhere. The hyperplanes
$w^\top u_n+b=0$ have pairwise nonproportional normals $(u_n,1)$.
At a point of the nth hyperplane on none of the others, crossing the
hyperplane produces the gradient jump $c_n(u_n,1)$; all other terms
are affine locally. The jump must vanish, so each $c_n=0$.

Let $\mathsf H_{ni}=h_i(u_n)$ be the unnormalized feature matrix.
The readout kernel is $(\mathsf H\mathsf H^\top/N)\otimes I_C$.
Gaussian ReLU fourth moments and Cauchy--Schwarz give
$\E\|\psi\|^4\le3N^2s_{\max}^4/2$. Independence of columns yields
\[
 \E\left\|\frac{\mathsf H\mathsf H^\top}{NM}-K_*\right\|_F^2
 \le\frac{3s_{\max}^4}{2M},\qquad
 \Pr\left(\left\|\frac{\mathsf H\mathsf H^\top}{NM}-K_*\right\|>
                  \kappa_*/2\right)
 \le\frac{6s_{\max}^4}{M\kappa_*^2}.
\]
For the encoder Jacobian, conditioning on the hidden weights kills the
cross terms in the independent centered readouts. Bounding each gate by
one gives $\E\|J_\thetap\|_F^2\le C_G$. Also
$\E\|e_0\|^2\le C_r$ by the activation second moment and
$\|z_0-y\|^2\le2\|z_0\|^2+2\|y\|^2$.
Markov and a three-event union bound prove the asserted event.
Apply Theorem~\ref{thm:astra_affine} on the full output space.
Finally $\operatorname{tr}(K_*)\le s_{\max}^2/2$ implies
$\kappa_*\le s_{\max}^2/(2N)$ and the threshold inequality.
\end{proof}
\begin{remark}
This proves a frozen-linearization instance only. Constants depend on
the fixed dataset and encoder; nearly duplicate bottleneck inputs can
make $\kappa_*$ arbitrarily small. Exact duplicates require a compatible
restricted-space formulation. The displayed constants are numerically
vacuous for the reported large datasets and are not empirical width thresholds.
\end{remark}

\begin{proposition}[Truncated-Chernoff alternative]
\label{prop:astra_chernoff}
For the same Gaussian/ReLU feature law, let $X_i=\psi_i\psi_i^\top/N$.
Choose $R_c>0$ such that
\[
 2N(R_c+2s_{\max}^2)e^{-R_c/(2s_{\max}^2)}\le\kappa_*/2.
\]
Then
\begin{equation}\label{eq:astra_chernoff}
 \Pr\left\{\lambda_{\min}\left(\sum_{i=1}^M X_i\right)
               <M\kappa_*/4\right\}
 \le N\exp[-M\kappa_*/(16R_c)].
\end{equation}
In particular $M\ge(16R_c/\kappa_*)\log(N/\eta)$ suffices for failure
probability $\eta\in(0,1)$. If used in the previous corollary or the
nonlinear theorem below, replace the head floor $\kappa_*M/2$ by
$\kappa_*M/4$ and propagate this changed constant.
\end{proposition}
\begin{proof}
Let $Z$ be the maximum of the N squared Gaussian preactivations.
Then $\|X_i\|\le Z$ and $\Pr(Z>u)\le2N e^{-u/(2s_{\max}^2)}$.
Integrating the tail gives
$\E[\|X_i\|1_{\{\|X_i\|>R_c\}}]
\le2N(R_c+2s_{\max}^2)e^{-R_c/(2s_{\max}^2)}$.
Thus $Y_i=X_i1_{\{\|X_i\|\le R_c\}}$ satisfies
$0\preceq Y_i\preceq R_cI$ and $\E Y_i\succeq\kappa_*I/2$.
Since each $X_i$ is positive semidefinite,
$\sum_iX_i\succeq\sum_iX_i1_{\{\|X_i\|\le R_c\}}$ and therefore
$\lambda_{\min}(\sum_iX_i)\ge\lambda_{\min}(\sum_iY_i)$; the
lower-tail bound for the truncated sum applies verbatim to the
untruncated sum. Apply matrix Chernoff \citep{astra_tropp2012} with
$\delta=1/2$ and $\mu_{\min}\ge M\kappa_*/2$.
\end{proof}
\begin{remark}
The constant is Tropp's $\delta=1/2$ lower tail after the standard
simplification
$[e^{-\delta}(1-\delta)^{-(1-\delta)}]^{\mu_{\min}/R_c}
\le e^{-\delta^2\mu_{\min}/(2R_c)}$.
Keeping the exact constant $-\log(e^{-1/2}\sqrt2)=0.15343\ldots$
replaces $1/16$ by approximately $1/13.04$; the difference is immaterial
at the widths involved. The one-sided ReLU tail also permits $N$ in
place of $2N$ in the truncation estimate. This certificate removes a
power of $\kappa_*^{-1}$ at the cost of logarithmic truncation factors;
it can still be numerically vacuous. Neither its threshold nor the
Frobenius threshold is a necessity theorem.
\end{remark}

\section{Interior curvature witnesses and normalization}
\label{sec:astra_interior}
\begin{lemma}[Pointwise interior witness; review Lemma 3.1]
\label{lem:astra_interior}
At N sites let $h_n\in\R^d$, $a_n=Wh_n+b\in\R^q$, and
$\hat y_n=\Phi_n(a_n)\in\R^C$, $C\ge2$.
Perturbing $a_n$ must not change logits at other sites in this lemma.
Fix a parameter point with finite logits and existing downstream
Jacobians $D_n=\partial\Phi_n/\partial a_n$. Let
$\Pi=I-\boldsymbol1\boldsymbol1^\top/C$,
$p_{\min,n}=\min_c p_{n,c}$. For a unit $a\in\R^q$ define
\[
 S_{h,a}=N^{-1}\sum_n p_{\min,n}\|\Pi D_na\|^2h_nh_n^\top,
 \qquad S_h^p=N^{-1}\sum_n p_{\min,n}h_nh_n^\top.
\]
If W belongs to $\phip$, then
\[
 \lambda_{\max}(G_{\phip\phip})\ge\lambda_{\max}(G_{WW})
 \ge\lambda_{\max}(S_{h,a}).
\]
If additionally $\|\Pi D_na\|\ge s$ for every n with $h_n\ne0$,
then the final lower bound is at least $s^2\lambda_{\max}(S_h^p)$.
\end{lemma}
\begin{proof}
For each unit $u\in\R^d$, $\Delta W=au^\top$ has unit Frobenius norm.
The CE Hessian satisfies $H_n\succeq p_{\min,n}\Pi$: indeed
$v^\top H_nv=\min_c\sum_jp_{n,j}(v_j-c)^2
\ge p_{\min,n}\min_c\sum_j(v_j-c)^2$.
Consequently the witness Rayleigh quotient is at least
$N^{-1}\sum_n(h_n^\top u)^2p_{\min,n}\|\Pi D_na\|^2
=u^\top S_{h,a}u$. Maximize over u. Zero-padding the W perturbation
proves the full-block inequality. The additional hypothesis gives
$S_{h,a}\succeq s^2S_h^p$.
\end{proof}
\begin{remark}[Mixing, masks, and nonsmooth points]
For a general interior perturbation, let $\mathcal R_h\operatorname{vec}
(\Delta W)$ be its full unnormalized preactivation stack, including all
sites that can affect the loss. If $D_\Phi$ is the full downstream
Jacobian and $P_{\rm out}$ selects the supervised logits, the valid
replacement is
\begin{equation}\label{eq:astra_full_witness}
 \lambda_{\max}(G_{\phip\phip})\ge
 \frac1{N_{\rm out}}\|H_\tau^{1/2}P_{\rm out}D_\Phi
       \mathcal R_h\operatorname{vec}(\Delta W)\|^2,
 \qquad\|\Delta W\|_F=1.
\end{equation}
This follows by the chain rule and the Rayleigh principle and includes
BN, convolutions, and ignored sites. A product lower bound requires a
restricted downstream gain on this same witness, with matching site
normalization; a smallest nonzero singular value does not protect a
witness in a nullspace. At ReLU kinks one must specify a gate selection;
for nondegenerate Gaussian preactivations the ordinary Jacobian exists
a.s. Sites with zero feature and zero bias contribute zero to the
incoming lower bound, but a Jacobian convention is still necessary.
\end{remark}

\begin{theorem}[Fixed bottleneck, growing interior width; review Theorem 3.2]
\label{thm:astra_interior_relu}
Fix integers $q\ge1$, $C\ge2$, and constants $R,a_0,s_w,s_v>0$,
$s_b\ge0$. For each M let deterministic features $h_n\in\R^M$ satisfy
$\|h_n\|^2\le R^2M$ and
$\lambda_{\max}(S_h)\ge a_0M$, where
$S_h=N^{-1}\sum_nh_nh_n^\top$.
Consider $\hat y_n=V\operatorname{ReLU}(Wh_n+b)$ with independent
$W_{ij}\sim\mathcal N(0,s_w^2/M)$, $b_i\sim\mathcal N(0,s_b^2)$,
$V_{ci}\sim\mathcal N(0,s_v^2/q)$, independent of the features.
At any identically zero preactivation fix a gate in $[0,1]$ consistently
in all Jacobian blocks; such sites have $h_n=0$ and cannot contribute
to the W witness below. Under averaged CE in these Euclidean
coordinates there exist $c_0,p_0>0$
independent of M,N with
\[
 \Pr\{\lambda_{\max}(G_{WW})\ge c_0a_0M\}\ge p_0,
 \qquad \E\lambda_{\max}(G_{\phip\phip})\ge p_0c_0a_0M
\]
whenever W belongs to $\phip$.
One explicit choice is
\[
 s_{\max}^2=s_w^2R^2+s_b^2,\quad L_0=\sqrt{48q s_{\max}^2},
 \quad c_0=e^{-4L_0}/(32C),\quad p_0=p_V/4,
\]
where $p_V=\Pr\{\|V\|_F\le2,\ \|\Pi Ve_1\|\ge1/2\}>0$.
These constants are \emph{existence-only and numerically vacuous}; this
is a positive-probability/expectation statement, not a probability
approaching one as M grows.
\end{theorem}
\begin{proof}
Fix a top unit eigenvector u of $S_h$, write
$a_n=(h_n^\top u)^2$, $A=N^{-1}\sum_na_n\ge a_0M$, and
$\chi_n=1_{\{(Wh_n+b)_1>0\}}$.
At sites with $a_n>0$ this preactivation is a nondegenerate centered
Gaussian. For $Q=N^{-1}\sum_na_n\chi_n$, $\E Q=A/2$ and $0\le Q\le A$,
hence $\Pr(Q\ge A/4)\ge1/3$ by splitting the expectation.
The exact moment bound
$\E\|\operatorname{ReLU}(Wh_n+b)\|^2\le qs_{\max}^2/2$
gives $\Pr(\|\operatorname{ReLU}(Wh_n+b)\|>L_0)\le1/96$.
Thus the weighted untame energy U has $\E U\le A/96$ and
$\Pr(U>A/8)\le1/12$. With probability at least $1/4$, active tame
sites therefore carry energy at least $A/8$. No independence across
sites was used.

Independently, the V event has positive probability: a small open ball
around a matrix with first column a unit class contrast and other
columns zero is contained in it. At tame sites on that event,
$\|\hat y_n\|_\infty\le2L_0$ and $p_{\min,n}\ge e^{-4L_0}/C$.
The unit witness $\Delta W=e_1u^\top$ has logit perturbation
$(h_n^\top u)\chi_nVe_1$, hence curvature at least
$(e^{-4L_0}/C)(1/4)(A/8)=c_0A$. Multiply the independent event
probabilities. Nonnegativity gives the expectation bound.
\end{proof}
\begin{remark}[Scope and a squared-loss counterpart]
A fixed channel direction may be inactive on every site, making its
bound zero on some draws. Adaptive directions can improve the witness,
but are not required for validity. For random upstream features, condition
on an event where both feature bounds hold; its probability multiplies
$p_0$. The theorem covers a pointwise ReLU bottleneck followed by a
linear readout, not an arbitrary convolutional BN decoder. It supplies
no residual-subspace coercivity for Theorem~\ref{thm:astra_affine}.
For scalar squared loss, independent readout entries
$v_i\sim\mathcal N(0,s_v^2/q)$ give the simpler bound
\[
 \E\lambda_{\max}(G_{WW})\ge\frac{s_v^2}{2q}\lambda_{\max}(S_h).
\]
Indeed, the same witness has curvature $v_1^2Q$, and independence gives
$\E(v_1^2Q)=(s_v^2/q)(A/2)$.
\end{remark}

\begin{proposition}[Counterexamples 3.3 and 3.4]
\label{prop:astra_bottleneck_counter}
A fixed ReLU bottleneck does not in general admit a positive interior
floor with probability $1-o_M(1)$. Moreover, train-mode BN can erase
an entire width-linear incoming witness even with a positive stabilizer.
\end{proposition}
\begin{proof}
For the first claim take all $h_n=h\ne0$ with $\|h\|^2=M$ and q
independent centered nondegenerate Gaussian preactivations. With
probability $2^{-q}$ all q gates are inactive at every site, and
$J_W=0$. This probability is independent of M; centered Gaussian bias
does not remove it.

For the second take $h_n=h$ throughout the \emph{whole normalization
group}, $\|h\|^2=M$, and place train-mode BN immediately after $Wh_n$.
For every positive stabilizer $\varepsilon>0$, its output is the BN
affine offset, independent of W. Thus every weight perturbation is
annihilated and $G_{WW}=0$, whereas
$S_h^p=p_{\min}hh^\top$ has eigenvalue $p_{\min}M$ if the downstream
logits have fixed finite nondegenerate probabilities. At
$\varepsilon=0$ this zero-variance forward pass is undefined; it is
not that counterexample. Constant inputs with a $1\times1$ convolution
or periodic boundaries give literal convolutional realizations.
\end{proof}

\begin{proposition}[BN derivative: an upper bound, not a lower witness]
\label{prop:astra_bn_derivative}
For one normalization group of size L, let
$P=I-\boldsymbol1\boldsymbol1^\top/L$, $t=Pu$,
$s^2=\|t\|^2/L$, and $\widetilde\sigma=\sqrt{s^2+\varepsilon}>0$,
with $\varepsilon\ge0$. With fixed affine scale $\gamma$ and offset,
\begin{equation}\label{eq:astra_bn_derivative}
 D_{\rm BN}=\frac{\gamma}{\widetilde\sigma}
 \left(P-\frac{tt^\top}{L(s^2+\varepsilon)}\right),\qquad
 \|D_{\rm BN}d\|^2\le\frac{\gamma^2}{s^2+\varepsilon}\|Pd\|^2.
\end{equation}
If $s^2>0$, decompose $Pd=d_\parallel+d_\perp$ parallel and orthogonal
to t. Then exactly
\[
 \|D_{\rm BN}d\|^2=\frac{\gamma^2}{s^2+\varepsilon}
 \left[\|d_\perp\|^2+
 \left(\frac{\varepsilon}{s^2+\varepsilon}\right)^2\|d_\parallel\|^2\right].
\]
\end{proposition}
\begin{proof}
Differentiate $\gamma Pu/\sqrt{\|Pu\|^2/L+\varepsilon}$.
The middle matrix annihilates constants, has eigenvalue one on
centered vectors orthogonal to t, and eigenvalue
$\varepsilon/(s^2+\varepsilon)$ in the t direction. This proves both
identities and the contraction upper bound. If $s^2=0$ and
$\varepsilon>0$, the derivative is simply $\gamma P/\sqrt\varepsilon$.
\end{proof}

\begin{proposition}[Function-preserving BN scale law]
\label{prop:astra_bn_scale}
Suppose a linear or convolutional block W is used downstream only
through its immediately following train-mode BN. Scale W and its bias
by c>0, keep the other parameters fixed, and co-scale the stabilizer
$\varepsilon\mapsto c^2\varepsilon$. On batches with defined forward
passes, logits are unchanged and
\begin{equation}\label{eq:astra_bn_scale}
 J_W(cW)=c^{-1}J_W(W),\qquad G_{WW}(cW)=c^{-2}G_{WW}(W),
\end{equation}
where the left side uses the co-scaled bias and stabilizer.
For $\varepsilon=0$ the same identities hold without a stabilizer change
provided every channel has positive variance.
\end{proposition}
\begin{proof}
Centering is linear and variance scales by $c^2$, so
$\operatorname{BN}_{c^2\varepsilon}(cu)=\operatorname{BN}_{\varepsilon}(u)$.
This identity holds as a function of W in a neighborhood where defined;
differentiate it and use unchanged output probabilities in the GGN.
The incoming bias Jacobian also scales by $c^{-1}$. Jacobians with
respect to parameters that are not co-scaled remain unchanged, by
differentiating the function identity with respect to those parameters.
\end{proof}
\begin{remark}[Known parameterization sensitivity]
The effective-learning-rate effect of normalization is established
\citep{vanlaarhoven2017l2}. The scalar
$\|W\|_F^2\lambda_{\max}(G_{WW})$ is invariant under this uniform
rescaling, not under arbitrary channelwise rescaling. For positive
diagonal D, channelwise scaling W to DW with channelwise stabilizers
co-scaled is also function preserving, but the GGN transforms by a
nonuniform inverse congruence; the scalar generally changes.
The audited four-channel example with scales $(0.2,1,3,8)$ changed it
by a factor 235 at identical logits. This is an example, not a universal
factor. At fixed positive $\varepsilon$, report each $s_i^2/\varepsilon$
and the scale-control departure: neither a universal accuracy threshold
such as $s_i^2\ge10^3\varepsilon$ nor a final-logit error bound follows
without bounds on c and on downstream sensitivity. Empirical tolerances
from a particular block are diagnostics only.
\end{remark}

\begin{proposition}[A constant centered Gram can coexist with growing curvature]
\label{prop:astra_center_counter}
For $M\ge2$, take three feature rows $(1,0)$, $(0,1)$, $(-1,-1)$ padded
with zeros to dimension M, and $w_M=e_1/\sqrt M$. Put a single train-mode
BN with $\gamma=1$, $\varepsilon=0$ after $w_M^\top h$, followed by
binary logits $(\operatorname{BN}(w_M^\top h),0)$. Then the centered
incoming Gram has top eigenvalue one, while
$\lambda_{\max}(G_{ww})=M c_*$ for a constant $c_*>0$.
\end{proposition}
\begin{proof}
The nonzero Gram block is
$\left(\begin{smallmatrix}2/3&1/3\\1/3&2/3\end{smallmatrix}\right)$.
At $w=e_1$ the centered preactivation is $(1,0,-1)$, with positive
variance, and the perturbation $e_2$ induces $(0,1,-1)$ with a
nonzero component transverse to it. The BN derivative and finite-logit
CE curvature give a positive Rayleigh quotient. Apply
Proposition~\ref{prop:astra_bn_scale} with $c=M^{-1/2}$.
This growth is precisely that function-preserving rescaling: the
normalization $\|w_M\|=M^{-1/2}$ is the whole mechanism. It refutes
an \emph{unqualified} necessity claim for a width-linear centered Gram,
not a claim that assumes width-uniform downstream gain. It is not a
normalization-free mechanism.
\end{proof}

\begin{proposition}[Conditional normalization-gain width law]
\label{prop:astra_bn_gain}
Let input dimension be d (for a fixed-size convolution, d is a fixed
multiple of $M+K$). On a fixed BN group let the \emph{centered} feature
covariance be $S_c$, and initialize a row w independently with
$\E w=0$, $\E ww^\top=\sigma_w^2 I_d/d$. Its pre-BN variance satisfies
\begin{equation}\label{eq:astra_bn_variance}
 \E_w s^2=\sigma_w^2\operatorname{tr}(S_c)/d.
\end{equation}
For a realized row assume $c_-/d\le s^2\le c_+/d$, with
$0<c_-\le c_+<\infty$, and
$0<\gamma_-\le|\gamma|\le\gamma_+$. Then
\begin{equation}\label{eq:astra_bn_gain}
 \frac{\gamma_-^2d}{c_++\varepsilon d}
 \le\frac{\gamma^2}{s^2+\varepsilon}
 \le\frac{\gamma_+^2d}{c_-+\varepsilon d}.
\end{equation}
For a given unit parameter witness whose pre-BN perturbation d' has
$\|d'_\perp\|^2\ge\tau\|d'\|^2$ with $\tau>0$, its post-BN energy
is at least
$\tau\gamma_-^2d\|d'\|^2/(c_++\varepsilon d)$.
A final GGN lower bound further requires a positive lower gain of the
actual downstream softmax-weighted, masked map on that post-BN perturbation.
\end{proposition}
\begin{proof}
$s^2=w^\top S_cw$, so taking a trace proves
\eqref{eq:astra_bn_variance}. Invert the two realized variance bounds
to obtain \eqref{eq:astra_bn_gain}. The transverse energy identity in
Proposition~\ref{prop:astra_bn_derivative} proves the witness assertion;
then apply \eqref{eq:astra_full_witness} with the stated downstream gain.
\end{proof}
\begin{remark}[Diagnostic model and its premises]
If $s^2=c_0/(M+K)$, its squared BN gain is exactly
\[
 \gamma^2(M+K)/(c_0+\varepsilon(M+K)).
\]
The premise is fixed positive \emph{total centered} incoming energy as
channels are added, together with control of the realized row variance.
Standard fan-in initialization with per-channel centered energy fixed
instead gives width-independent \emph{expected} pre-BN variance.
A fixed total uncentered second moment is insufficient (constant
features have zero centered variance). An expectation of variance cannot
be inverted into a concentration statement about gain. Even with fixed
pre-BN witness energy, a radial witness is annihilated at
$\varepsilon=0$ and need not inherit the gain. These qualifications
make this a conditional diagnostic identity, not a new general BN theorem
or an unconditional replacement for Lemma~\ref{lem:phi_scaling}.
\end{remark}
\section{Residual excitation and phase-restricted attribution}
\label{sec:astra_excitation}
\begin{lemma}[Directional gradient gain is not directional GGN curvature]
\label{lem:astra_direction_gain}
Let $R_u$ isometrically embed a directional parameter subspace and
$J_u=J_\thetap R_u$. Then
\[
 \|J_u^\top r\|^2\le\lambda_{\max}(J_u^\top J_u)\|r\|^2.
\]
If $H\succeq0$ and $r\in\operatorname{range}(H)$, also
\begin{equation}\label{eq:astra_pseudoinverse}
 \|J_u^\top r\|^2\le\lambda_{\max}(J_u^\top HJ_u)r^\top H^\dagger r.
\end{equation}
For finite-logit softmax CE the range condition holds. A bound
$H\succeq h_0\Pi_{\rm cls}$ on the class-contrast space, $h_0>0$,
permits $r^\top H^\dagger r\le\|r\|^2/h_0$.
\end{lemma}
\begin{proof}
The first claim is the operator-norm inequality. For the second, write
\[J_u^\top r=(H^{1/2}J_u)^\top H^{\dagger/2}r.\]
Apply the operator-norm inequality again.
The last assertion follows by diagonalizing H on its range.
Here the Hessian floor is understood for the softmax Hessian, whose
range is exactly the class-contrast space at finite logits.
\end{proof}
\begin{remark}[CE counterexample and saturation]
For one example with logits $(\theta,0)$, $\theta=0$, and label one, the residual and Jacobian are
\[r=(-1/2,1/2),\qquad J=(1,0)^\top.\]
Then
$|\nabla_\theta\loss|^2=1/4$, while
$G_{\theta\theta}\|r\|^2=(1/4)(1/2)=1/8$.
Thus replacing the logit Gram by the GGN without the factor in
\eqref{eq:astra_pseudoinverse} is false. At general first-class probability
p the ratio is $1/(2p(1-p))$, unbounded near zero or one.
A uniform positive h0 is a further assumption, not supplied by
saturation. For a linear encoder and unit input direction u,
$R_u b=\operatorname{vec}(bu^\top)$, so $J_u^\top r=(\nabla_\Theta\loss)u$.
\end{remark}

\begin{theorem}[Cumulative energy and excitation; review Theorem 2.1]
\label{thm:astra_excitation}
Let $A_u,A_v$ be constant directional logit Jacobians along a path,
$r(t)=\nabla_z\ell(z(t))$ square integrable on $[0,T]$, and
$C_u=A_uA_u^\top$, $\lambda_u=\|C_u\|$, with analogous v notation.
Define $Q_T=\int_0^T r(t)r(t)^\top dt$ and
$E_u=\int_0^T\|A_u^\top r(t)\|^2dt$. Then
\[
 E_u=\operatorname{tr}(C_uQ_T)\le\lambda_u\operatorname{tr}(Q_T).
\]
Suppose $\lambda_v>0$, $\lambda_u>0$, $\operatorname{tr}(Q_T)>0$,
and a top unit eigenvector $q_v$ of $C_v$ satisfies
$q_v^\top Q_Tq_v\ge\rho\operatorname{tr}(Q_T)$ with $\rho>0$.
Then $E_v\ge\rho(\lambda_v/\lambda_u)E_u$; division by $E_u$ requires
$E_u>0$. More generally, if $C_v\succeq c\lambda_vP$ and
$\operatorname{tr}(PQ_T)\ge\rho\operatorname{tr}(Q_T)$ for a projector P,
replace $\rho$ by $c\rho$.
\end{theorem}
\begin{proof}
Integrate $\|A_u^\top r\|^2=r^\top C_ur$ and use
$C_u\preceq\lambda_uI$. Since $C_v\succeq\lambda_vq_vq_v^\top$,
$E_v\ge\lambda_v q_v^\top Q_Tq_v\ge\rho\lambda_v\operatorname{tr}(Q_T)$.
The projector version is identical.
\end{proof}
\begin{remark}
This theorem is loss-agnostic under its constant-Jacobian and integrability
hypotheses; it includes affine CE. An envelope
$\|r(t)\|\le C_r/(1+\mu t)$ gives $E_u\le\lambda_u C_r^2/\mu$.
Here $\lambda_u$ is a logit-kernel eigenvalue. The exponential residual
formulas in the next counterexamples specifically use squared loss.
Excitation is an output-space condition on the visited residual; an
input principal component alone does not establish it.
\end{remark}

\begin{proposition}[Counterexamples 2.2 and 2.3 and the decoupled ratio]
\label{prop:astra_energy_counter}
For affine squared loss take
$C_v=L e_1e_1^\top$, $C_u=\ell e_2e_2^\top$, $L,\ell>0$, with
total kernel $\operatorname{diag}(L,\ell)$.
If $r(0)=e_2$, then $E_v(T)=0<E_u(T)$ for $T>0$.
If $r(0)=e_1+e_2$, then
\[
 E_v(T)=(1-e^{-2LT})/2,\qquad E_u(T)=(1-e^{-2\ell T})/2,
\]
so their ratio tends to one even for arbitrarily large $L/\ell$.
More generally assume e1,e2 are orthonormal eigenvectors of the total
constant kernel with positive eigenvalues $\gamma_v,\gamma_u$,
$C_v=\lambda_v e_1e_1^\top$, $C_u=\lambda_u e_2e_2^\top$, and the
initial residual components are $a_v,a_u$ with $a_u\ne0$,
$\lambda_u>0$, $T>0$. Then
\begin{equation}\label{eq:astra_decoupled}
 \frac{E_v(T)}{E_u(T)}=
 \frac{\lambda_v}{\lambda_u}\frac{a_v^2}{a_u^2}
 \frac{\gamma_u}{\gamma_v}
 \frac{1-e^{-2\gamma_vT}}{1-e^{-2\gamma_uT}}.
\end{equation}
\end{proposition}
\begin{proof}
Each residual component is $a_je^{-\gamma_jt}$; integrate
$\lambda_ja_j^2e^{-2\gamma_jt}$. The first two constructions are
realized by the affine map $(\sqrt L\theta_1,\sqrt\ell\theta_2)$.
Without the decoupling hypothesis, these scalar exponentials do not
represent mixed residual components.
\end{proof}

\begin{proposition}[Attribution during a residual-excited phase]
\label{prop:astra_phase}
Assume the differentiability, absolute continuity, unit-flow and
energy-identity conventions of Corollary~\ref{cor:astra_attribution},
but not its global Rayleigh bounds. Let
$q_b(t)=r^\top K_br$, $D(t)=q_\thetap(t)+q_\phip(t)$, and
$\Delta L=\int_0^TD(t)dt=\loss(0)-\loss(T)>0$.
For a measurable phase $E\subset[0,T]$, put
\[
 \tau_E=|E|,\qquad w_E=\frac{\int_E D(t)dt}{\Delta L}.
\]
On E suppose a measurable top-k projector $P_k(t)$ of $K_\phip(t)$
(with a fixed measurable tie convention) satisfies
\[
 \|P_k(t)r(t)\|^2\ge\alpha\|r(t)\|^2,\quad
 \lambda_k(K_\phip(t))\ge\ell_*>0,\quad
 q_\thetap(t)\le a\|r(t)\|^2,
\]
where $0<\alpha\le1$, $a\ge0$. Then, when the phase denominator is positive,
\begin{align}
 \frac{\int_E q_\thetap}{\int_E D}&\le\frac{a}{a+\alpha\ell_*},
 \label{eq:astra_phase}\\
 \frac{\int_0^Tq_\thetap}{\Delta L}
 &\le w_E\frac{a}{a+\alpha\ell_*}+(1-w_E).
 \label{eq:astra_phase_total}
\end{align}
A verified complementary share bound b replaces the last term by
$(1-w_E)b$. Time-varying pointwise constants give the sharper
$D(t)dt/\Delta L$-weighted average of their share caps.
\end{proposition}
\begin{proof}
PSD spectral decomposition gives
$q_\phip\ge\lambda_k\|P_kr\|^2\ge\alpha\ell_*\|r\|^2$.
Thus $q_\thetap\le[a/(a+\alpha\ell_*)]D$ on E.
Integrate on E and its complement, where $q_\thetap\le D$.
The phase duration is explicit but imposes no value on $w_E$.
For $E=[0,\tau]$, $\int_E D=\loss(0)-\loss(\tau)$; under squared
loss additionally $\|r(t)\|\le e^{-\alpha\ell_*t}\|r(0)\|$
through this interval by differentiating its squared norm.
\end{proof}
\begin{corollary}[What a width-dependent cumulative share additionally requires]
\label{cor:astra_phase_width}
For a family satisfying Proposition~\ref{prop:astra_phase}, suppose
its early-phase constants obey $a\le a_*$ and
$\alpha\ell_*\ge cM$, with fixed $a_*\ge0,c>0$. On the complement
assume $0<b_-\le q_\thetap/D\le b_+\le1$ at almost every time with
$D>0$, for constants independent of M. Writing the cumulative share
as $S_M$, one has
\[
 b_-(1-w_E)\le S_M
 \le w_E\frac{a_*}{a_*+cM}+b_+(1-w_E).
\]
Consequently $S_M=O(M^{-1})$ holds if and only if
$1-w_E=O(M^{-1})$. The same statement holds for nonnegative weighted
sums on one discrete sequence, interpreting $w_E$ and $S_M$ as
\emph{gradient-energy} fractions.
\end{corollary}
\begin{proof}
The upper bound is Proposition~\ref{prop:astra_phase}. For the lower
bound retain only the complementary numerator and use its share floor.
Both implications follow since $b_->0$ and $a_*/(a_*+cM)=O(M^{-1})$.
This characterizes what the claimed cumulative scaling needs under the
stated late-phase floor; early excitation or a short phase alone is insufficient.
\end{proof}

\begin{remark}[Discrete diagnostic and admissible reference space]
Replacing integrals by nonnegative weighted sums of gradient squares
proves the same \emph{gradient-energy} inequalities on any one sampled
sequence, regardless of optimizer. This is not Adam loss-decrease
attribution. Probe-defined phases and training-minibatch energies are
different sequences unless explicitly matched. No phase-share bound
alone replaces the stability assumptions needed to compare two trajectories.
For CE let $\Pi_{\rm cls}=I_N\otimes(I_C-\boldsymbol1\boldsymbol1^\top/C)$.
The residual belongs to its rank-$N(C-1)$ range. For a uniformly random
unit vector in that range, expected energy in a projector P is
$\operatorname{tr}(P\Pi_{\rm cls})/[N(C-1)]$, since its second moment is
$\Pi_{\rm cls}/[N(C-1)]$. The full-space baseline
$\operatorname{rank}(P)/(NC)$ is generally inappropriate.
\end{remark}

\section{A solvable serial cross-entropy instance}
\label{sec:astra_serial}
\begin{proposition}[Information retained by a fitting serial encoder]
\label{prop:astra_information}
If the predictions depend on X only through $Z=f_\thetap(X)$ and a
classifier $q_\phip(Y\mid Z)$ has population CE at most $\eta<\infty$,
then $H(Y\mid Z)\le\eta$ and $I(Y;Z)\ge H(Y)-\eta$.
The same statement holds for the empirical distribution with a fixed model.
\end{proposition}
\begin{proof}
CE is conditional entropy plus the nonnegative expected conditional KL
divergence. Suppression of a selected coordinate or failure under a
specified shift is therefore different from absence of all training-label
information in the serial representation.
\end{proof}

\begin{theorem}[Serial CE instance; review Theorem 5.1]
\label{thm:astra_serial}
Let $Y$ be uniform on $\{-1,1\}$ and let training inputs be
$x_1=(S,0)$, $x_2=(0,C_{\rm ctx})$ with $S=C_{\rm ctx}=Y$.
The shared linear encoder is $W=(a,v)$, giving $z_1=aS,z_2=vC_{\rm ctx}$.
For integer $M\ge1$ define
\[
 F=z_1+\sum_{j=1}^M\beta_jz_2=aS+bvC_{\rm ctx},\quad
 b=\sum_j\beta_j,\quad\thetap=(a,v),\quad\phip=\beta.
\]
All coordinates follow unit Euclidean gradient flow on the averaged
binary CE $\loss=\E\log(1+e^{-YF})=\log(1+e^{-q})$, $q=a+bv$,
from $a_0=0,v_0=1,\beta_0=0$, with no clipping, decay or momentum.
Fix $0<\delta<1/2$, $m=\log((1-\delta)/\delta)$, and let $T_m$ be
the first time q reaches m. Write $\varrho=(1+e^q)^{-1}$ and
$\Psi(q)=q+e^q-1$. Then:
\begin{align}
 \dot a&=\varrho,&\dot v&=b\varrho,&\dot b&=Mv\varrho,\label{eq:astra_serial_ode}\\
 v&=\cosh(\sqrt M a),&b&=\sqrt M\sinh(\sqrt M a),\label{eq:astra_serial_phase}\\
 q&=a+\frac{\sqrt M}{2}\sinh(2\sqrt M a),&
 \dot q&=(M+1+2b^2)\varrho.\label{eq:astra_serial_q}
\end{align}
The solution is global and reaches m in finite time, with
\begin{equation}\label{eq:astra_serial_time}
 \frac{\Psi(m)}{M+1+2m^2}\le T_m\le\frac{\Psi(m)}{M+1}.
\end{equation}
At that time, writing $a_M=a(T_m)$,
\begin{gather}
 0<a_M\le\frac{m}{M+1},\qquad b(T_m)v(T_m)\ge\frac{Mm}{M+1},
 \qquad Ma_M\longrightarrow m,\label{eq:astra_serial_update}\\
 0\le v(T_m)-1\le\frac{m^2}{2M},\qquad
 \|W(T_m)-W_0\|\le
 \sqrt{\frac{m^2}{(M+1)^2}+\frac{m^4}{4M^2}}.\label{eq:astra_serial_disp}
\end{gather}
For fixed m, $a_M$ decreases strictly with M; removing the contextual
path entirely requires $a=m$ at the same fitting margin. Moreover,
\begin{equation}\label{eq:astra_serial_envelope}
 \varrho(t)\le\frac{1/2}{1+(M+1)t/4},\qquad t\ge0.
\end{equation}
In the stated parameter blocks and loss convention,
\[
 \lambda_{\max}(G_{\thetap\thetap})=\varrho(1-\varrho)(1+b^2),\quad
 \lambda_{\max}(G_{\phip\phip})=M\varrho(1-\varrho)v^2,
\]
so at initialization they equal $1/4,M/4$ and the disparity is M.
For the separately trained frozen encoder $W=(0,1)$, $\beta_F(0)=0$,
its margin obeys $\Psi(q_F(t))=Mt$ and
\begin{equation}\label{eq:astra_serial_gap}
 0\le q(t)-q_F(t)\le\frac{(1+2m^2)t}{2+Mt/2}
 \le\frac{2(1+2m^2)}M,\qquad0\le t\le T_m.
\end{equation}
Under the specified reversal $S=Y,C_{\rm ctx}=-Y$, the joint model at
$T_m$ has error one; the fitted spectral-only comparator has error zero.
\end{theorem}
\begin{proof}
The training margin is the same on each example, so direct differentiation
gives \eqref{eq:astra_serial_ode}. Divide by $\dot a=\varrho>0$ to get
$dv/da=b$, $db/da=Mv$, which proves \eqref{eq:astra_serial_phase}.
It implies $b^2=M(v^2-1)$ and \eqref{eq:astra_serial_q}.
Since $0<\dot a\le1/2$, a stays finite on every finite time interval;
the phase formula then gives global existence. The strictly increasing
map $Q_M(a)=a+(\sqrt M/2)\sinh(2\sqrt M a)$ maps $[0,\infty)$
onto itself and has a finite inverse at m. On this compact a interval,
$\dot a\ge\delta>0$, so the hitting time is finite.

$\sinh x\ge x$ gives $Q_M(a)\ge(M+1)a$, strictly for $a>0$.
Before fitting, $v\ge1$, $0\le bv=q-a\le m$, hence $b\le m$.
The invariant gives $v-1=b^2/[M(v+1)]\le m^2/(2M)$.
For fixed positive a the contextual term increases strictly with M;
invert to get strict decrease of $a_M$. Its upper bound implies
$\sqrt M a_M\to0$, and the expansion
$m=(M+1)a_M+O(M^2a_M^3)$ gives $Ma_M\to m$.

$\dot\varrho=-(M+1+2b^2)\varrho^2(1-\varrho)$ and
$\varrho\le1/2$ imply $(d/dt)(1/\varrho)\ge(M+1)/2$.
Integrate from $\varrho(0)=1/2$ for the envelope.
Since $\Psi'(q)=1/\varrho$, before fitting
$\frac d{dt}\Psi(q)=M+1+2b^2\in[M+1,M+1+2m^2]$;
integration proves \eqref{eq:astra_serial_time}.
The margin gradients are $Y(1,b)$ and $Yv\boldsymbol1_M$;
the scalar logistic Hessian is $\varrho(1-\varrho)$, giving the GGN.

For the frozen system $\dot q_F=M/(1+e^{q_F})$. Thus
$0\le\Psi(q)-\Psi(q_F)=\int_0^t(1+2b^2)ds\le(1+2m^2)t$.
Convexity gives a lower bound $\Psi'(q_F)(q-q_F)$ for this difference.
As $q_F\le e^{q_F}-1$,
$Mt=\Psi(q_F)\le2(e^{q_F}-1)$, so
$\Psi'(q_F)\ge2+Mt/2$, proving \eqref{eq:astra_serial_gap}.
The shifted margin is $a-bv=2a_M-m<0$: for $M>1$ use
\eqref{eq:astra_serial_update}, and for $M=1$ use the strict
inequality $Q_1(a)>2a$. The spectral-only shifted margin is m.
\end{proof}

\begin{corollary}[Analytic instantiation of the displacement bound]
\label{cor:astra_serial_bound}
For Theorem~\ref{thm:astra_serial}, embed the scalar prediction in
symmetric logits $(F/2,-F/2)$. The packet residual norm is
$R=\sqrt2\varrho$, and on $[0,T_m]$
\[
 \|J_\thetap\|=\sqrt{(1+b^2)/2}\le B_m:=\sqrt{(1+m^2)/2}.
\]
With envelope amplitude $C_r=1$ and $\mu=(M+1)/4$,
Theorem~\ref{thm:twoscale} gives
\begin{equation}\label{eq:astra_serial_bound}
 \|W(T_m)-W_0\|\le\frac{4B_m}{M+1}\log\frac1{\sqrt2\delta}.
\end{equation}
Also $\|W(T_m)-W_0\|\le\sqrt{1+m^2}\,a_M
\le\sqrt{1+m^2}m/(M+1)$.
\end{corollary}
\begin{proof}
The symmetric logit derivative has squared norm $(1+b^2)/2$, while
the full residual norm is $\sqrt2\varrho$, so their product is the
exact encoder gradient norm. The envelope has true initial amplitude
$1/\sqrt2$; $C_r=1$ is a valid larger amplitude matching the main
text's $C_r\ge1$ convention. Its target is $\sqrt2\delta<1$.
All constants apply on the fitting region $|b|\le m$.
Alternatively integrate $\|\dot W\|=\varrho\sqrt{1+b^2}
\le\sqrt{1+m^2}\dot a$.
\end{proof}
\begin{remark}[Embedding, head, and initialization conventions]
For logits $(F,0)$, loss, residual norm and GGN agree, but the unprojected
logit-Jacobian norm is $\sqrt{1+b^2}$; the symmetric embedding makes
the norm-product bound exact. A freely trained dense two-row head is
not the parameter block $\phip=\beta$: its initial head eigenvalue is
$M/2$, not $M/4$, while the encoder eigenvalue remains $1/4$ at the
same initial function. Labels need not be balanced for the training
ODE; balance is used in the information/noisy-accuracy statement below.
The initial encoder orientation in Theorem~\ref{thm:astra_serial} is
explicitly asymmetric; the next corollary removes that choice for W,
not the routing/readout asymmetry.
\end{remark}

\begin{proposition}[Readout normalization and representation interpretation]
\label{prop:astra_serial_control}
Replace the contextual readout by $\alpha_M\sum_j\gamma_jz_2$,
$\gamma_j(0)=0$, at unit rates. Its reduced equation for
$b=\alpha_M\sum_j\gamma_j$ is $\dot b=M\alpha_M^2v\varrho$.
In particular $\alpha_M=M^{-1/2}$ gives exactly the $M=1$ dynamics
for every M. At every finite M in the original theorem, the noiseless
spectral coordinate $a_MS$ still determines Y. If it is instead
observed as $Z_S=a_MY+\sigma\xi$, with $\xi\sim\mathcal N(0,1)$
independent, $\sigma>0$ fixed, then Bayes accuracy is
$\Phi(a_M/\sigma)$ and $I(Y;Z_S)\le a_M^2/(2\sigma^2)$.
\end{proposition}
\begin{proof}
Differentiate each readout and sum. The Gaussian likelihood-ratio test
is thresholding at zero, giving the accuracy. For
$Q=\mathcal N(0,\sigma^2)$,
$\E_Y\operatorname{KL}(P_{Z_S|Y}\|Q)=a_M^2/(2\sigma^2)
=I(Y;Z_S)+\operatorname{KL}(P_{Z_S}\|Q)$.
This noise is a specified observation perturbation, not training noise.
\end{proof}

\begin{corollary}[General and isotropic encoder initialization]
\label{cor:astra_isotropic}
In the serial model let $m>0$, $a_0<m$, $v_0\ne0$, and allow any
readout initialization with $\sum_j\beta_{j0}=0$. Set
$u=a-a_0$, $\Delta=m-a_0>0$. Then
\begin{align}
 v&=v_0\cosh(\sqrt M u),\quad b=\sqrt Mv_0\sinh(\sqrt M u),
 \label{eq:astra_iso_phase}\\
 q&=a_0+u+\frac{\sqrt Mv_0^2}{2}\sinh(2\sqrt M u),\notag\\
 0<u_M&\le\frac{\Delta}{1+Mv_0^2},\qquad
 Mu_M\longrightarrow\frac{\Delta}{v_0^2},\qquad
 \frac{a(T_m)-a_0}{m-a_0}\longrightarrow0.
 \label{eq:astra_iso_update}
\end{align}
The fitting root decreases strictly with M and
\begin{equation}\label{eq:astra_iso_disp}
 |v(T_m)-v_0|\le\frac{\Delta^2}{2M|v_0|^3},\qquad
 \|W(T_m)-W_0\|\le
 \sqrt{\frac{\Delta^2}{(1+Mv_0^2)^2}+
       \frac{\Delta^4}{4M^2|v_0|^6}}.
\end{equation}
With $\Psi_{a_0}(q)=q-a_0+e^q-e^{a_0}$,
\[
 \frac{\Psi_{a_0}(m)}{1+Mv_0^2+2\Delta^2/v_0^2}
 \le T_m\le\frac{\Psi_{a_0}(m)}{1+Mv_0^2}.
\]
For a frozen encoder at $(a_0,v_0)$ and the same initial readouts,
$\Psi_{a_0}(q_F)=Mv_0^2t$. If
$p_0=e^{a_0}/(1+e^{a_0})$ and $C_0=1+2\Delta^2/v_0^2$, then
\begin{equation}\label{eq:astra_iso_gap}
 0\le q_J-q_F\le\frac{C_0t}{1+e^{a_0}+p_0Mv_0^2t}
 \le\frac{C_0}{p_0Mv_0^2},\qquad0\le t\le T_m.
\end{equation}
Under contextual reversal the fitted margin is $2a(T_m)-m$.
For $a_0<m/2$, error is one for all sufficiently large M, with sufficient
condition $Mv_0^2>m/(m-2a_0)$. For $m/2\le a_0<m$, the reversed
classifier is correct at every width. If $a_0\ge m$, define fitting to
stop immediately, in which case it is also correct.

If $(a_0,v_0)\sim\mathcal N(0,\sigma^2 I_2)$, $\sigma>0$, then
$v_0\ne0$ almost surely; use the preceding stopping convention.
The expected reversal error satisfies
\begin{equation}\label{eq:astra_iso_error}
 \lim_{M\to\infty}\E_{W_0}\operatorname{Err}_{\rm reversal}
 =\Phi\left(\frac{m}{2\sigma}\right).
\end{equation}
The same-threshold spectral-only comparator has zero reversal error.
\end{corollary}
\begin{proof}
Dividing the ODE by $\dot a>0$ gives
\eqref{eq:astra_iso_phase}; only the initial readout sum enters.
Each $\beta_i-\beta_j$ is conserved, so equal initial readouts are not
needed for closure. Positivity of the contextual margin
$b v=(\sqrt Mv_0^2/2)\sinh(2\sqrt M u)$ and the sinh inequality
prove \eqref{eq:astra_iso_update}; Taylor expansion at
$\sqrt M u_M\to0$ gives the limit. The invariant is
$b^2=M(v^2-v_0^2)$; v keeps its sign and $|v|\ge|v_0|$.
Since $bv\le\Delta$, $|b|\le\Delta/|v_0|$ and
\[
 |v-v_0|=\frac{b^2}{M(|v|+|v_0|)}
 \le\frac{\Delta^2}{2M|v_0|^3}.
\]
This finite-M inequality proves \eqref{eq:astra_iso_disp}.
Now $\dot q=(1+Mv_0^2+2b^2)\varrho$ and
$b^2\le\Delta^2/v_0^2$ on the fitting interval; integrate the
transformed margin for the time bounds. Global existence follows as
before from $\dot u\le1$ and the phase formula; every finite m is hit.

The transformed joint/frozen difference lies between zero and $C_0t$.
Also $q_F-a_0\le(e^{q_F}-e^{a_0})/e^{a_0}$ gives
$\Psi'_{a_0}(q_F)\ge1+e^{a_0}+p_0Mv_0^2t$.
Convexity proves \eqref{eq:astra_iso_gap}. For $a_0<m/2$, use
$a(T_m)\le a_0+\Delta/(1+Mv_0^2)$ to get the finite-M reversal
condition. For $a_0>m/2$ the classifier is correct at every fitted
width, since $u_M>0$ gives $2a(T_m)-m>2a_0-m>0$; equality
$a_0=m/2$ also gives a positive margin. Immediate stopping covers
$a_0\ge m$. Almost surely the limiting error is
$1_{\{a_0<m/2\}}$, apart from the zero-probability boundary and
$v_0=0$ events. Dominated convergence gives \eqref{eq:astra_iso_error}.
\end{proof}

\begin{remark}[Confidence dependence and limiting quantity]
For $a_0\ne0$, $a(T_m)/a_0\to1$ but $a(T_m)/m\to a_0/m$;
the vanishing quantity is the \emph{update} fraction in
\eqref{eq:astra_iso_update}. For $0<\rho<1$, a Gaussian tail and its
maximum density give, with probability at least $1-\rho$,
\[
 |a_0|\le\sigma\sqrt{2\log(4/\rho)},\qquad
 |v_0|\ge\frac{\rho\sigma}{2}\sqrt{\pi/2}.
\]
Substitute these bounds in the finite-M inequalities for uniform
confidence bounds. The dependence is severe: for fixed
$(\Delta,v_0)\in(0,\infty)\times(\R\setminus\{0\})$, expansion
of the root yields
\[
 Mu_M=\frac{\Delta}{v_0^2}\left[
 1-\frac1{Mv_0^2}-\frac{2\Delta^2}{3Mv_0^4}+O(M^{-2})\right].
\]
Indeed substitute $u=A/M+B/M^2$ in the sinh equation and match its
constant and $M^{-1}$ coefficients. The expansion's remainder is
pointwise, not already uniform in rho. Its small-$|v_0|$ crossover
requires control of $\Delta^2/(Mv_0^4)$; at the confidence cutoff this
has scale $\Delta^2\rho^{-4}\sigma^{-4}/M$, a fourth power of
$1/\rho$, not a second. A rigorous uniform remainder can instead be
bounded from the sinh series on this controlled argument. Do not
present this asymptotic crossover as a universal necessary width.
\end{remark}

\paragraph{Provenance and interpretation.}
Architecture- and initialization-dependent implicit bias is established
in linear-network analyses \citep{astra_yun2021,astra_moroshko2020};
trajectory-level incremental learning is studied by
\citet{astra_berthier2023}. Exact deep-linear dynamics
\citep{astra_saxe2014} and conservation of layer-norm differences
\citep{astra_du2018balance} supply the underlying tools. CE feature
competition is central to \citet{pezeshki2021gradient}.
The serial example instantiates a known optimization-metric mechanism
with finite-margin formulas, an encoder-update comparison and a specified
reversal distribution. Replication is equivalent to a rate M on b,
and normalization removes it; the function class does not grow.
These particular calculations are not a claim of a new general implicit-bias
mechanism, and finite-time or serial composition alone is not a novelty claim.

\section{A scoped nonlinear existence result}
\label{sec:astra_nonlinear}

\paragraph{Flow convention and normalization.}
Here the loss is squared loss, the encoder is linear, the decoder has
one hidden ReLU layer, and every parameter uses unit Euclidean gradient
flow. There is no BatchNorm, convolution, weight decay, clipping,
preconditioning, or stochastic sampling. Write
\[
 f_n=V\sigma(W\Theta x_n+b)+\beta,\qquad
 \loss=\frac1{2N}\sum_{n=1}^N\norm{f_n-y_n}^2,
 \qquad \thetap=\Theta,\quad\phip=(W,b,V,\beta).
\]
Parameter norms are Euclidean, with Frobenius norms on matrices.
Set $\bar f=N^{-1/2}(f_1,\ldots,f_N)$ and
$e=\bar f-\bar y$, so that $\loss=\norm e^2/2$.
At a ReLU kink choose a measurable gate in $[0,1]$, using the same gate
in every parameter block for that neuron and datum. A solution is an
absolutely continuous parameter path satisfying the resulting gradient
flow almost everywhere. The theorem concerns every such solution;
uniqueness is not asserted. Local existence and the energy identity
needed below are proved in the proof. This formulation avoids identifying
arbitrary backpropagation selections with the Clarke subdifferential.

\begin{theorem}[Theorem N: nonlinear, unnormalized-readout existence result]
\label{thm:astra_nonlinear}
Fix positive integers $N,D,K,C$, fixed data $x_n\in\R^D,y_n\in\R^C$,
and a deterministic $\Theta_0\in\R^{K\times D}$ such that the
$\Theta_0x_n$ are pairwise distinct. Fix $s_w,s_v,s_b>0$ and initialize
independently across neurons, entries, and blocks:
\[
 w_i\sim\mathcal N(0,s_w^2 I_K/K),\quad
 b_i\sim\mathcal N(0,s_b^2),\quad
 V_{ci}\sim\mathcal N(0,s_v^2/M).
\]
The output bias $\beta_0$ is fixed, or independent of the other blocks
with an $M$-independent distribution of finite second moment. In
particular $\beta_0$ need not be centered. Compare joint training from
this initialization with training of all decoder parameters from the
same initialization while fixing $\Theta=\Theta_0$.
For every $\rho\in(0,1)$ there are finite positive constants
$\kappa,A,B,R_0,C_{\rm gap},M_*$, depending on the fixed data,
initialization scales and $\rho$, but not on $M$, with the following
property. For each $M\ge M_*$, with probability at least $1-\rho$,
both flows have global solutions, and every pair of solutions obeys,
for every $t\ge0$,
\begin{align*}
 \norm{e^J(t)},\norm{e^F(t)}
   &\le R_0e^{-\kappa Mt/2}, \tag{N1}\label{eq:astra_n1}\\
 \sup_t\norm{\Theta^J(t)-\Theta_0}
   &\le\frac{2AR_0}{\kappa M},&
 \sup_t\norm{\phip^{J,F}(t)-\phip_0}
   &\le\frac{2BR_0}{\kappa\sqrt M}, \tag{N2}\label{eq:astra_n2}\\
 \sup_t\norm{\bar f^J(t)-\bar f^F(t)}
   &\le\frac{C_{\rm gap}R_0}{\kappa M}. \tag{N3}\label{eq:astra_n3}
\end{align*}
For every finite $T>0$ with $\loss(0)>\loss(T)$ on the joint path,
\begin{equation*}
 \frac{\int_0^T\norm{\nabla_{\thetap}\loss}^2dt}
 {\loss(0)-\loss(T)}
 \le\frac{A^2}{A^2+\kappa M/2}. \tag{N4}\label{eq:astra_n4}
\end{equation*}
For $0<\delta<R_0$, residual norm at most $\delta$ is reached by
$2(\kappa M)^{-1}\log(R_0/\delta)$. If the initial residual is already
at most $\delta$, the first hitting time can be zero.
The constants and threshold constructed below are existence-only and
can be numerically vacuous. This is not a uniform theorem for learned
encoders, growing datasets, general deep decoders, or production Adam.
\end{theorem}

\begin{proof}
We give explicit finite constants to make all quantifiers and
probability events checkable.

\emph{Population kernel and initialization events.}
Let $X=\max_n\norm{x_n}$, $z_n=\Theta_0x_n$, and
\[
 H_{ni}=N^{-1/2}\sigma(w_i^Tz_n+b_i),\quad
 K_* =\E[H_{:i}H_{:i}^T],\quad
 \kappa_* =\lambda_{\min}(K_*),\quad \kappa=\kappa_*/2.
\]
Here $K_*$ uses one random neuron with the displayed normalization,
so that $\E[HH^T]=MK_*$. The distinct-input kink argument in the proof
of Corollary~\ref{cor:astra_init} gives $\kappa_*>0$: if a linear
combination of the ReLU features vanishes almost surely, it vanishes
everywhere by continuity and full Gaussian support in $(w,b)$;
one can cross each distinct affine hyperplane away from the others,
forcing its coefficient to vanish.
Put
\[
 s_{\max}^2=\frac{s_w^2}{K}\max_n\norm{z_n}^2+s_b^2,
 \quad Z_0=\sqrt{1+(\norm{\Theta_0}X)^2},\quad Z=Z_0+X.
\]
Operator norm may be used for $\norm{\Theta_0}$ here; its Frobenius
norm is also a valid upper bound. For $\norm{\Theta-\Theta_0}\le1$,
\begin{equation}\label{eq:astra_zbound}
 \sqrt{1+\norm{\Theta x_n}^2}\le Z.
\end{equation}
Indeed, with $a=\norm{\Theta_0}X$,
$(\sqrt{1+a^2}+X)^2-[1+(a+X)^2]
=2X(\sqrt{1+a^2}-a)\ge0$.

Choose positive $h_0,C_w,C_v,A_0,R_0$ such that each of
\begin{equation}\label{eq:astra_fiveevents}
 \norm{H_0}_F\le h_0\sqrt M,\quad
 \norm{W_0}_F\le C_w\sqrt M,\quad
 \norm{V_0}_F\le C_v,\quad
 \norm{J_{\Theta,0}}_{\rm op}\le A_0,\quad
 \norm{e_0}\le R_0
\end{equation}
has failure probability at most $\rho/7$. Here and below $J$ is
the Jacobian of $\bar f$. For example, Markov's inequality applies
with respective second-moment upper bounds
\[
 \frac{s_{\max}^2}{2},\quad s_w^2,\quad Cs_v^2,\quad
 Cs_v^2s_w^2X^2,\quad
 Cs_v^2s_{\max}^2+2\E\norm{\beta_0}^2+2\norm{\bar y}^2
\]
after dividing the first two squared norms by $M$. Take each squared
constant at least $7/\rho$ times the corresponding bound, enlarging
to a positive value if necessary. For $J_{\Theta,0}$, conditional
centering and independence of the entries of $V_0$ eliminate cross
terms in $V_0\Gamma_{n,0}W_0$; its expected squared Frobenius norm
is at most $Cs_v^2s_w^2$. The factors $x_n$ and $N^{-1/2}$ give the
stated bound. The residual bound follows by conditioning on $(W_0,b_0)$
and using the independent centered $V_0$; the output bias itself need
not be centered.
The variance calculation used in Corollary~\ref{cor:astra_init} gives
\[
 \Pr\{\lambda_{\min}(H_0H_0^T)<\kappa M\}
 \le\frac{6s_{\max}^4}{M\kappa_*^2}.
\]
Thus $M\ge42s_{\max}^4/(\rho\kappa_*^2)$ makes this sixth failure
probability at most $\rho/7$.

\emph{Boundary event and constants.}
Set
\[
 h=h_0+1,\quad \Omega=Z(C_v+1),\quad
 B=\sqrt{h^2+\Omega^2+1},\quad R_\phi=4BR_0/\kappa,
 \quad \tau_i=\frac{X\norm{w_{i,0}}+ZR_\phi}{\sqrt M}.
\]
For the seventh event define
\[
 \mathcal B=\sum_{i=1}^M\norm{V_{:i,0}}\norm{w_{i,0}}
 \mathbf1\{\min_n|w_{i,0}^Tz_n+b_{i,0}|\le\tau_i\}.
\]
Conditioning on $w_{i,0}$, the Gaussian density of $b_{i,0}$ is bounded
by $(s_b\sqrt{2\pi})^{-1}$. A union bound over $n$, independence
of $V_{:i,0}$, and $\E\norm{V_{:i,0}}\le s_v\sqrt{C/M}$ give
\[
 \E\mathcal B\le C_{\rm bd}
 :=\frac{2Ns_v\sqrt C}{s_b\sqrt{2\pi}}
 \left(X\E\norm w^2+ZR_\phi\E\norm w\right).
\]
Here $\E\norm w^2=s_w^2$ and $\E\norm w\le s_w$; either the actual
moment or its bound can define the constant. Put
$\mathcal B_*=7C_{\rm bd}/\rho$. Then
$\Pr\{\mathcal B>\mathcal B_*\}\le\rho/7$.
Define
\begin{align*}
 A&=A_0+X\{R_\phi C_w+(C_v+1)+\mathcal B_*\},&
 R_\theta&=4AR_0/\kappa,\\
 D_H&=XC_wR_\theta+ZR_\phi,&
 D_K&=(2h_0+1)D_H,\\
 C_{\rm gap}&=2D_K+2\Omega^2+A^2.
\end{align*}
It suffices to take the integer ceiling of
\begin{equation}\label{eq:astra_widththreshold}
 M_*\ge\max\left\{1,R_\theta^2,R_\phi^2,D_H,
              2D_K/\kappa,42s_{\max}^4/(\rho\kappa_*^2)\right\}.
\end{equation}
All constants are determined before $M$ is chosen. The intersection
of the five events in \eqref{eq:astra_fiveevents}, the head-kernel
event, and the boundary event has probability at least $1-\rho$.
Initialization at a kink has probability zero and may be excluded
without another probability charge.

\emph{Deterministic estimates in the bootstrap tube.}
Work on this intersection and in the closed tube
\[
 \norm{\Theta-\Theta_0}\le R_\theta/M,\qquad
 \norm{\phip-\phip_0}\le R_\phi/\sqrt M.
\]
The width conditions imply $\norm{\Theta-\Theta_0}\le1$ and
$\norm{\phip-\phip_0}\le1$, as well as
$R_\theta/M\le M^{-1/2}$. Write $\omega=(W,b)$.
ReLU is 1-Lipschitz, so
\begin{align*}
 \norm{H-H_0}_F
 &\le X\norm{W_0}_F\norm{\Theta-\Theta_0}
        +Z\norm{\omega-\omega_0}
 \le D_H/\sqrt M,\\
 \norm H_F&\le h\sqrt M,\qquad
 \norm{HH^T-H_0H_0^T}_{\rm op}\le D_K.
\end{align*}
For the second line use $D_H\le M$ and expand the product; in
particular $2h_0D_H+D_H^2/M\le D_K$.
Consequently
\begin{equation}\label{eq:astra_headfloor}
 HH^T\succeq(\kappa M/2)I_N.
\end{equation}
The preactivation change for neuron $i$ on any datum is bounded by
\[X\norm{w_{i,0}}R_\theta/M+Z\norm{\omega_i-\omega_{i,0}}\le\tau_i.\]
Thus a changed gate belongs to the initial boundary event above.
To control the encoder Jacobian use the exact three-term decomposition
\begin{equation}\label{eq:astra_three_terms}
 V\Gamma_nW-V_0\Gamma_{n,0}W_0
 =(V-V_0)\Gamma_nW_0+V\Gamma_n(W-W_0)
       +V_0(\Gamma_n-\Gamma_{n,0})W_0.
\end{equation}
The last term has operator norm at most $\mathcal B$. The first
uses the current gates and initial $W_0$; the second uses the current
$V$; both distinctions matter. Averaging over $n$ and multiplying by
$x_n$ yields
\[
 \norm{J_\Theta}_{\rm op}
 \le A_0+X\{\norm{V-V_0}_F\norm{W_0}_F
       +\norm V_F\norm{W-W_0}_F+\mathcal B_*\}\le A.
\]
The separate block bounds are
\[
 \norm{J_\omega}_{\rm op}\le Z\norm V_F\le\Omega,\qquad
 \norm{J_V}_{\rm op}\le h\sqrt M,\qquad
 \norm{J_\beta}_{\rm op}=1,
 \quad \norm{J_\phip}_{\rm op}\le B\sqrt M.
\]
Furthermore $J_\phip J_\phip^T\succeq
(HH^T)\otimes I_C\succeq(\kappa M/2)I_{NC}$.
These estimates also hold for the frozen path in its decoder tube.

\emph{Flow existence, chain rule and continuation.}
For completeness, local solutions in the convention above can be
obtained by replacing ReLU by
$\sigma_\epsilon(u)=(u+\sqrt{u^2+\epsilon^2})/2$.
For fixed $M$, on a compact parameter neighborhood its values and
first derivatives are uniformly bounded for $0<\epsilon\le1$.
The smooth gradient vector fields are therefore uniformly bounded
there, giving a common positive interval of existence and uniformly
Lipschitz paths on a smaller neighborhood. A subsequence converges
uniformly. The finitely many gates have a weak-* convergent subsequence
in $L^\infty$ with limits in $[0,1]$ and the correct values away from
zero preactivations. Each block gradient is linear in the corresponding
gate, with all remaining coefficients converging uniformly. Passing
to the integral equations gives one common limiting gate selection
and a local solution.
Along any such absolutely continuous solution, each preactivation
has derivative zero almost everywhere on its zero level set.
Consequently $d\sigma(g(t))/dt=\gamma(t)\dot g(t)$ almost everywhere,
for every admissible gate selection. The ordinary product rule then
gives $\dot e=-JJ^Te$ and
$\dot\loss=-\norm{J^Te}^2$ almost everywhere. These facts follow
directly here; related nonsmooth chain-rule frameworks are discussed
by \citet{astra_davis2020,astra_bolte2021}. We do not infer a Clarke
inclusion solely from the computational graph.

Until a first exit from the tube,
$\frac d{dt}\norm e^2\le-\kappa M\norm e^2$, proving (N1).
The integrated speeds are bounded by
\[
 \int_0^\infty\norm{\dot\Theta}\,dt\le\frac{2AR_0}{\kappa M}
     =\frac{R_\theta}{2M},\qquad
 \int_0^\infty\norm{\dot\phip}\,dt
     \le\frac{2BR_0}{\kappa\sqrt M}
     =\frac{R_\phi}{2\sqrt M}.
\]
At a first exit these bounds contradict the tube boundary. Thus no
exit occurs. Bounded velocities in this compact tube preclude finite
escape, give a limit at any finite endpoint, and the local existence
argument continues the path from that endpoint. This proves global
existence and (N1)--(N2) for every selection, without uniqueness.

\emph{Common linear reference and attribution.}
Let
\[
 K_0=(H_0H_0^T)\otimes I_C+
             (\mathbf1\mathbf1^T/N)\otimes I_C,
 \qquad e^{\rm lin}(t)=e^{-K_0t}e_0.
\]
The second summand is the constant output-bias kernel.
For the full output kernels $K^J(t)$ and $K^F(t)$,
\[
 \norm{K^J(t)-K_0}_{\rm op}\le D_K+\Omega^2+A^2,
 \qquad
 \norm{K^F(t)-K_0}_{\rm op}\le D_K+\Omega^2.
\]
Since $K_0\succeq\kappa MI$, variation of constants gives
\[
 e^a(t)-e^{\rm lin}(t)
 =-\int_0^t e^{-K_0(t-s)}(K^a(s)-K_0)e^a(s)\,ds,
 \qquad a\in\{J,F\}.
\]
Use $\norm{e^a(s)}\le R_0$ and
$\int_0^te^{-\kappa M(t-s)}ds\le1/(\kappa M)$ and add the two
bounds to obtain (N3). Both paths track the same linear reference;
(N3) is not a general stability theorem for arbitrary nonlinear
joint and frozen flows. Applying
Corollary~\ref{cor:astra_attribution} with $a=A^2$ and decoder floor
$\kappa M/2$ gives (N4), completing the proof.
\end{proof}

\begin{remark}[Threshold and scope]
The threshold \eqref{eq:astra_widththreshold} is deliberately
conservative. The boundary term contains an explicit factor $N/\rho$
and $R_0$ can scale as $\rho^{-1/2}$, before its further dependence
on $R_\phi$, the data, and $\kappa_*$ is substituted. When these
factors are nonzero and the remaining scales are held fixed, already
$R_\theta^2$ can contain a contribution of order
$N^2/(\rho^3\kappa_*^2)$; the full displayed construction can be
substantially worse. This is an illustration of vacuity, not a
universal necessary lower bound (degenerate inputs can remove some
terms), and it is not an empirical width prediction.
\end{remark}

\paragraph{Relation to prior work and a normalization control.}
Positive limiting kernels, small parameter movement, and control of
ReLU gate changes are established tools in wide-network convergence
proofs \citep{astra_du2019,astra_arora2019}. Linearized dynamics are
studied by \citet{astra_lee2019}; their dependence on parameterization
and the distinction between lazy and feature-learning limits are
central to \citet{chizat2019lazy,yang2021tensor}.
The divergent kernel scale of the naive standard parameterization and
its relation to layerwise training rates are also explicit in
\citet{astra_sohldickstein2020}.
The result above is a scoped instantiation with a coupled finite-width
encoder and a joint/frozen comparison. It does not establish a new
general lazy-training mechanism. Its gain disparity vanishes under
an explicit $1/\sqrt M$ readout in the following precise convention:
write $f=M^{-1/2}U\sigma(W\Theta x+b)+\beta$, initialize the entries
of $U$ with variance $s_v^2$, and train $U$ at unit Euclidean rate.
This preserves the initial function distribution but changes the
optimization metric. The readout kernel is then $HH^T/M=O(1)$,
and the encoder kernel remains $O(1)$; there is no forced factor-$M$
readout/encoder separation. This does not assert equality of the two
kernels or of their directional gains. Merely inserting the factor
while retaining variance $s_v^2/M$ for the trained readout would be
a different, shrinking-output initialization.

% BEGIN ASTRA BIBTEX: copy the entries below, without the leading percent, to references.bib.
% @inproceedings{astra_karimi2016,
%   author = {Karimi, Hamed and Nutini, Julie and Schmidt, Mark},
%   title = {Linear Convergence of Gradient and Proximal-Gradient Methods Under the {Polyak--Lojasiewicz} Condition},
%   booktitle = {European Conference on Machine Learning and Knowledge Discovery in Databases},
%   year = {2016}, url = {https://arxiv.org/abs/1608.04636}
% }
% @article{astra_tropp2012,
%   author = {Tropp, Joel A.},
%   title = {User-Friendly Tail Bounds for Sums of Random Matrices},
%   journal = {Foundations of Computational Mathematics}, volume = {12},
%   pages = {389--434}, year = {2012}, doi = {10.1007/s10208-011-9099-z}
% }
% @article{vanlaarhoven2017l2,
%   author = {van Laarhoven, Twan},
%   title = {{L2} Regularization versus Batch and Weight Normalization},
%   journal = {arXiv preprint arXiv:1706.05350}, year = {2017},
%   url = {https://arxiv.org/abs/1706.05350}
% }
% @inproceedings{astra_yun2021,
%   author = {Yun, Chulhee and Krishnan, Shankar and Mobahi, Hossein},
%   title = {A Unifying View on Implicit Bias in Training Linear Neural Networks},
%   booktitle = {International Conference on Learning Representations}, year = {2021},
%   url = {https://arxiv.org/abs/2010.02501}
% }
% @inproceedings{astra_moroshko2020,
%   author = {Moroshko, Edward and Woodworth, Blake and Gunasekar, Suriya and Lee, Jason D. and Srebro, Nathan and Soudry, Daniel},
%   title = {Implicit Bias in Deep Linear Classification: Initialization Scale vs Training Accuracy},
%   booktitle = {Advances in Neural Information Processing Systems}, volume = {33}, year = {2020},
%   url = {https://arxiv.org/abs/2007.06738}
% }
% @article{astra_berthier2023,
%   author = {Berthier, Rapha{\"e}l},
%   title = {Incremental Learning in Diagonal Linear Networks},
%   journal = {Journal of Machine Learning Research}, volume = {24}, number = {171},
%   pages = {1--26}, year = {2023}, url = {https://jmlr.org/papers/v24/22-1395.html}
% }
% @inproceedings{astra_saxe2014,
%   author = {Saxe, Andrew M. and McClelland, James L. and Ganguli, Surya},
%   title = {Exact Solutions to the Nonlinear Dynamics of Learning in Deep Linear Neural Networks},
%   booktitle = {International Conference on Learning Representations}, year = {2014},
%   url = {https://arxiv.org/abs/1312.6120}
% }
% @inproceedings{astra_du2018balance,
%   author = {Du, Simon S. and Hu, Wei and Lee, Jason D.},
%   title = {Algorithmic Regularization in Learning Deep Homogeneous Models: Layers are Automatically Balanced},
%   booktitle = {Advances in Neural Information Processing Systems}, volume = {31}, year = {2018},
%   url = {https://arxiv.org/abs/1806.00900}
% }
% @inproceedings{pezeshki2021gradient,
%   author = {Pezeshki, Mohammad and Kaba, Sekou-Oumar and Bengio, Yoshua and Courville, Aaron and Precup, Doina and Lajoie, Guillaume},
%   title = {Gradient Starvation: A Learning Proclivity in Neural Networks},
%   booktitle = {Advances in Neural Information Processing Systems}, volume = {34}, year = {2021},
%   url = {https://arxiv.org/abs/2011.09468}
% }
% @inproceedings{astra_du2019,
%   author = {Du, Simon S. and Zhai, Xiyu and Poczos, Barnabas and Singh, Aarti},
%   title = {Gradient Descent Provably Optimizes Over-parameterized Neural Networks},
%   booktitle = {International Conference on Learning Representations}, year = {2019},
%   url = {https://openreview.net/forum?id=S1eK3i09YQ}
% }
% @article{astra_davis2020,
%   author = {Davis, Damek and Drusvyatskiy, Dmitriy and Kakade, Sham and Lee, Jason D.},
%   title = {Stochastic Subgradient Method Converges on Tame Functions},
%   journal = {Foundations of Computational Mathematics}, volume = {20}, pages = {119--154},
%   year = {2020}, url = {https://arxiv.org/abs/1804.07795}
% }
% @article{astra_bolte2021,
%   author = {Bolte, J{\'e}r{\^o}me and Pauwels, Edouard},
%   title = {Conservative Set Valued Fields, Automatic Differentiation, Stochastic Gradient Methods and Deep Learning},
%   journal = {Mathematical Programming}, volume = {188}, pages = {19--51},
%   year = {2021}, doi = {10.1007/s10107-020-01501-5},
%   url = {https://arxiv.org/abs/1909.10300}
% }
% @inproceedings{astra_lee2019,
%   author = {Lee, Jaehoon and Xiao, Lechao and Schoenholz, Samuel S. and Bahri, Yasaman and Novak, Roman and Sohl-Dickstein, Jascha and Pennington, Jeffrey},
%   title = {Wide Neural Networks of Any Depth Evolve as Linear Models Under Gradient Descent},
%   booktitle = {Advances in Neural Information Processing Systems}, volume = {32}, year = {2019},
%   url = {https://arxiv.org/abs/1902.06720}
% }
% @inproceedings{astra_arora2019,
%   author = {Arora, Sanjeev and Du, Simon S. and Hu, Wei and Li, Zhiyuan and Wang, Ruosong},
%   title = {Fine-Grained Analysis of Optimization and Generalization for Overparameterized Two-Layer Neural Networks},
%   booktitle = {Proceedings of the 36th International Conference on Machine Learning},
%   series = {Proceedings of Machine Learning Research}, volume = {97}, pages = {322--332}, year = {2019},
%   url = {https://proceedings.mlr.press/v97/arora19a.html}
% }
% @inproceedings{chizat2019lazy,
%   author = {Chizat, Lenaic and Oyallon, Edouard and Bach, Francis},
%   title = {On Lazy Training in Differentiable Programming},
%   booktitle = {Advances in Neural Information Processing Systems}, volume = {32}, year = {2019},
%   url = {https://arxiv.org/abs/1812.07956}
% }
% @inproceedings{yang2021tensor,
%   author = {Yang, Greg and Hu, Edward J.},
%   title = {Tensor Programs {IV}: Feature Learning in Infinite-Width Neural Networks},
%   booktitle = {Proceedings of the 38th International Conference on Machine Learning},
%   series = {Proceedings of Machine Learning Research}, volume = {139}, pages = {11727--11737}, year = {2021},
%   url = {https://proceedings.mlr.press/v139/yang21c.html}
% }
% @article{astra_sohldickstein2020,
%   author = {Sohl-Dickstein, Jascha and Novak, Roman and Schoenholz, Samuel S. and Lee, Jaehoon},
%   title = {On the Infinite Width Limit of Neural Networks with a Standard Parameterization},
%   journal = {arXiv preprint arXiv:2001.07301}, year = {2020},
%   url = {https://arxiv.org/abs/2001.07301}
% }
% END ASTRA BIBTEX
